%%%%%%%%%%%%%%%%%%%%%%%%%%%%%%%%%%%%%%%%%%%%%%%%%%%%%%%%%%%%%%%%%%%%%%
% Lớp 11 | Dãy số. Cấp số cộng. Cấp số nhân | Cấp số nhân
% Dạng 1D2?3-6: Tính tổng của cấp số nhân
% Số câu: 100  (sinh tự động bởi tach_cau.py; từ nay thầy sửa trực tiếp file này)
%%%%%%%%%%%%%%%%%%%%%%%%%%%%%%%%%%%%%%%%%%%%%%%%%%%%%%%%%%%%%%%%%%%%%%

%%%%%%%%%%%%%%%% TRẮC NGHIỆM 4 PHƯƠNG ÁN %%%%%%%%%%%%%%%%

% nguồn: Ok_Toan11_TracNghiem_Lan2
\begin{ex}%[1D2N3-6]%[ma:7efe0af6]
	Cho cấp số nhân $\left(u_n\right)$ có $u_1=-3$ và $q=-2$. Tính  tổng $10$ số hạng đầu tiên của cấp số nhân.
	 \choice
		 {$S_{10}=-511$}
		{\True  $S_{10}=1\,023$}
		{$S_{10}=1\,025$} 
		{$S_{10}=-1\,025$} 
	 \loigiai{
		Ta  có $S_{10}=u_1\cdot \dfrac{1-q^n}{1-q}=-3\cdot \dfrac{1-\left(-2\right)^{10}}{1-\left(-2\right)}=1\,023$.
	 }
\end{ex}

% nguồn: Ok_Toan11_TracNghiem_Lan2
\begin{ex}%[1D2N3-6]%[ma:a9bec4c7]
	Cho cấp số nhân có các số hạng lần lượt là $\dfrac{1}{4}$; $\dfrac{1}{2}$; $1$; $\cdots$; $2\,048$. Tính tổng $S$ của tất  cả các số hạng của cấp số nhân đã cho.
	\choice 
		 {\True $S=2\,047{,}75$}
		 {$S=2\,049{,}75$}
		{$S=4\,095{,}75$} 
		 {$S=4\,096{,}75$}
	\loigiai{ 
		Cấp số  nhân đã cho có $\heva{&u_1=\dfrac{1}{4}\\&q=2.}$\\
		 Suy ra $2\,048=2^{11}=u_1q^{n-1}=\dfrac{1}{2}{2^{n-1}}=2^{n-2}\Leftrightarrow n=13$.\\
		Hay cấp  số nhân đã cho có tất cả $13$ số hạng.\\
		Vậy $S_{13}=u_1\dfrac{1-q^{13}}{1-q}=\dfrac{1}{4}\cdot \dfrac{1-2^{13}}{1-2}=2\,047{,}75$. 
	 }
\end{ex}

% nguồn: Ok_TracNghiemLop11_Lan1
\begin{ex}%[1D2N3-6]%[ma:0a42d818]
	Cho cấp số nhân $(u_n)$ có số hạng đầu $u_1=3$ và công bội $q=2$. Tổng của $10$ số hạng đầu tiên  của cấp số nhân đó bằng
	 \choice
		 {$\dfrac{1023}{2}$}
		{$1023$} 
		 {$1536$}
		 {\True $3069$}
	 \loigiai{
		Ta có $S_{10}=u_1\cdot  \dfrac{1-q^{10}}{1-q}=3\cdot \dfrac{1-2^{10}}{1-2}=3069$.
	 }
\end{ex}

% nguồn: Ok_TracNghiemLop11_Lan1
\begin{ex}%[1D2N3-6]%[ma:834b8b4e]
	Cho cấp số nhân $\left(u_n\right)$ với công  bội $q\neq1$. Đặt $S_n=u_1+u_2+\cdots+u_n$. Khẳng định nào sau đây đúng?
	\choice 
		 {\True $S_n=\dfrac{u_1\left(1-q^n\right)}{1-q}$}
		{$S_n=\dfrac{u_1\left(1-q^{n-1}\right)}{1-q}$} 
		{$S_n=\dfrac{u_1(1-q)}{1-q^n}$} 
		{$S_n=u_1\left(1-q^n\right)$} 
	 \loigiai{
		Giả sử  $\left(u_n\right)$ là một cấp số nhân có công bội $q\neq1$. Đặt $S_n=u_1+u_2+\cdots+u_n$, khi đó
		 \[S_n=\dfrac{u_1\left(1-q^n\right)}{1-q}.\]
	} 
\end{ex}

% nguồn: Ok_TracNghiemLop11_Lan1
\begin{ex}%[1D2N3-6]%[ma:d1797b57]
	Cho cấp số nhân $\left( u_n\right)$ có  số hạng đầu $u_1=6$ và công bội $q=\dfrac{1}{2}$. Tính tổng của $8$ số hạng đầu tiên của cấp số nhân đã cho.
	\choice 
		 { $-12$}
		{  \True $\dfrac{765}{64}$}
		 { $- \dfrac{3}{64}$}
		 { $\dfrac{19}{2}$}
	 \loigiai{
		Ta có  $S_8=\dfrac{u_1\left( 1-q^8\right) }{1-q}=\dfrac{6\left(1-\left( \frac{1}{2}\right)^8 \right) }{1- \frac{1}{2}}=\dfrac{765}{64}$.
	 }
\end{ex}

% nguồn: Ok_Toan11_TracNghiem_Lan2
\begin{ex}%[1D2H3-6]%[ma:7ad41357]
	Cho cấp số nhân $\left(u_n\right)$ có $u_3 = 12$, $u_5 = 48$ và có công bội âm.  Tổng $7$ số hạng đầu của cấp số nhân đã cho bằng
	\choice 
		{\True $129$} 
		{$-129$} 
		{$128$} 
		{$-128$} 
	 \loigiai{
		Ta có: $u_4^2 = u_3 \cdot  u_5 = 576$. Vì $u_3 > 0$, $u_5 > 0$ và công bội âm nên: $u_4 = -24 \Rightarrow q = -2$.\\
		Lại có:  $u_3 = u_1 q^2 \Rightarrow u_1 = \dfrac{u_3}{q^2} = \dfrac{12}{4} = 3$.\\
		Áp dụng công thức  ta có: $S_7 = u_1\dfrac{1 - q^7}{1 - q} = 3 \cdot \dfrac{1 - (-2)^7}{1 - (-2)} = 129$.
	} 
\end{ex}

% nguồn: Ok_Toan11_TracNghiem_Lan2
\begin{ex}%[1D2H3-6]%[ma:f4ff9d1b]
	Số thập phân vô hạn tuần hoàn  $3{,}1555\ldots = 3{,}1(5)$ viết dưới dạng số hữu tỉ là:
	 \choice
		{$\dfrac{63}{20}$} 
		{\True  $\dfrac{142}{45}$}
		 {$\dfrac{1}{18}$}
		 {$\dfrac{7}{2}$}
	 \loigiai{
		$3{,}1555\ldots = 3{,}1 + 0{,}05 + 0{,}005 + 0{,}0005  + \ldots$ \\
		Dãy số $0{,}05; 0{,}005; 0{,}0005; \ldots$ là một cấp số nhân lùi vô hạn có $u_1 = 0{,}05$, $q = 0{,}1$.\\ 
		 Vậy $3{,}1555\ldots = 3{,}1 + \dfrac{0{,}05}{1 - 0{,}1} = \dfrac{142}{45}$.
	 }
\end{ex}

% nguồn: Ok_Toan11_TracNghiem_Lan2
\begin{ex}%[1D2H3-6]%[ma:7b0acb23]
	Cho cấp số nhân có các số hạng lần lượt là $1$, $4$, $16$, $64$, $\ldots$ Gọi $S_n$ là tổng của  $n$ số hạng đầu tiên của cấp số nhân đó. Mệnh đề nào sau đây đúng?
	 \choice
		{\True  $S_n = \dfrac{4^n - 1}{3}$}
		{$S_n = 4^{n -  1}$}
		{$S_n = \dfrac{n\left(1 +  4^{n - 1}\right)}{2}$}
		{$S_n = \dfrac{4\left(4^n  - 1\right)}{3}$}
	 \loigiai{
		Cấp số nhân đã cho có $\heva{&u_1 = 1\\ &q = 4} \Rightarrow S_n = u_1 \cdot \dfrac{1 - q^n}{1  - q} = 1 \cdot \dfrac{1 - 4^n}{1 - 4} = \dfrac{4^n - 1}{3}$.
	 }
\end{ex}

% nguồn: Ok_Toan11_TracNghiem_Lan2
\begin{ex}%[1D2H3-6]%[ma:316a1424]
	Tổng vô hạn sau đây $S=2+\dfrac{2}{3}+\dfrac{2}{3^2}+\cdots+\dfrac{2}{3^n}+\cdots$ có giá trị bằng mấy? 
	 \choice
		{$2$} 
		{$4$} 
		{$\dfrac{8}{3}$} 
		{\True $3$} 
	\loigiai{Ta có $S=2+\dfrac{2}{3}+\dfrac{2}{3^2}+\cdots+\dfrac{2}{3^n}+\cdots=2.\lim\dfrac{1-\dfrac{1}{3^n}}{1-\dfrac{1}{3}}=2\cdot\dfrac{1}{\dfrac{2}{3}}=2\cdot\dfrac{3}{2}=3$. 
	 }
\end{ex}

% nguồn: Ok_Toan11_TracNghiem_Lan2
\begin{ex}%[1D2H3-6]%[ma:b5052d79]
	Cho $\left(u_n\right)$ là cấp số nhân, đặt $S_n=u_1+u_2+\cdots +u_n$. Biết $S_2=4;\,S_3=13$  và $u_2<0$, giá trị $S_5$ bằng
	 \choice
		{$2$} 
		{\True $\dfrac{181}{16}$} 
		 {$\dfrac{35}{16}$}
		{$121$} 
	 \loigiai{
		Gọi $u_1$, $q$  lần lượt là số hạng đầu tiên và công bội của cấp số nhân cần tìm.\\
		Từ giả thiết ta  có $\heva{&S_2=4\\&S_3=13}\Leftrightarrow \heva{&u_1(1+q)=4\\&u_1\left(1+q+q^2\right)=13}\Leftrightarrow \heva{&u_1(1+q)=4\\&\hoac{&q=3\\&q=-\dfrac{3}{4}.}}$\\
		 Vì $\heva{&u_2<0\\&u_3=S_3-S_2=9>0} \Rightarrow q=\dfrac{u_3}{u_2}<0$ nên cấp số nhân cần tìm có $\heva{&u_1=16\\&q=-\dfrac{3}{4}.}$\\
	Do  đó $S_5=u_1\left(\dfrac{1-q^5}{1-q}\right)=\dfrac{181}{16}$.}
\end{ex}

% nguồn: Ok_Toan11_TracNghiem_Lan2
\begin{ex}%[1D2H3-6]%[ma:8529e96a]
	Giá trị của tổng $S=1+3+3^2+\cdots +3^{2\,018}$ bằng 
	 \choice
		 {\True $S=\dfrac{3^{2\,019}-1}{2}$}
		{$S=\dfrac{3^{2\,018}-1}{2}$} 
		{$S=\dfrac{3^{2\,020}-1}{2}$} 
		 {$S=-\dfrac{3^{2\,018}-1}{2}$}
	\loigiai{ 
		 Ta thấy $S$ là tổng của $2\,019$ số hạng đầu tiên của cấp số nhân với số hạng đầu là $u_1=1$, công bội $q=3$.\\
		Áp dụng công thức tính  tổng của cấp số nhân ta có $S=1\cdot \dfrac{1-3^{2\,019}}{1-3}=\dfrac{3^{2\,019}-1}{2}$.
	} 
\end{ex}

% nguồn: Ok_Toan11_TracNghiem_Lan2
\begin{ex}%[1D2H3-6]%[ma:7c6051e7]
	Cho cấp số  nhân $(u_n)$ thỏa mãn $\heva{&u_1 + u_2 + u_3 = 13\\ &u_4 - u_1 = 26}$. Tổng $8$ số hạng đầu của cấp số nhân $(u_n)$ là
	\choice 
		{$S_8  = 1\,093$}
		 {$S_8 = 3\,820$}
		{$S_8  = 9\,841$}
		 {\True $S_8 = 3\,280$}
	\loigiai{ 
		Ta có $\heva{&u_1 + u_2 + u_3 = 13\\ &u_4 - u_1 = 26} \Leftrightarrow \heva{&u_1 + u_1.q + u_1.q^2 = 13\\ &u_1.q^3 - u_1 = 26} \Leftrightarrow \heva{&u_1\left(1 + q + q^2\right) = 13\\ &u_1(q - 1)\left(1 + q  + q^2\right) = 26}$\\
		$\Leftrightarrow \heva{&u_1\left(1  + q + q^2\right) = 13\\ &q = 3} \Leftrightarrow \heva{&u_1 = 1\\ &q = 3}$.\\
		Vậy $S_8  = \dfrac{u_1\left(1 - q^8\right)}{1 - q} = \dfrac{1\left(1 - 3^8\right)}{1 - 3} = 3\,280$.
	} 
\end{ex}

% nguồn: Ok_Toan11_TracNghiem_Lan2
\begin{ex}%[1D2H3-6]%[ma:04c28d6d]
	Biết rằng $S = 1 + 2 \cdot 3 + 3 \cdot 3^2 + \ldots + 11 \cdot 3^{10} = a +  \dfrac{21 \cdot 3^b}{4}$. Tính $P = a + \dfrac{b}{4}$.
	 \choice
		 {$P = 1$}
		{$P = 2$} 
		{\True $P = 3$} 
		{$P  = 4$}
	 \loigiai{
		Từ giả thiết suy ra $3S  = 3 + 2 \cdot 3^2 + 3 \cdot 3^3 + \ldots + 11 \cdot 3^{11}$.\\
		 Do đó:
		\[ 
		-2S = S - 3S  = 1 + 3 + 3^2 + \ldots + 3^{10} - 11 \cdot 3^{11} = \dfrac{1 - 3^{11}}{1 - 3} - 11 \cdot 3^{11} =-\dfrac{1}{2} -\dfrac{1}{2} \cdot 21 \cdot 3^{11}.
		\] 
		Suy ra $S = \dfrac{1}{4}  + \dfrac{21 \cdot 3^{11}}{4}$.\\
		Vì $S = \dfrac{1}{4} + \dfrac{21 \cdot 3^{11}}{4} = a + \dfrac{21 \cdot 3^b}{4}$ nên $a = \dfrac{1}{4}$, $b  = 11$. \\
		Vậy $P = \dfrac{1}{4}  + \dfrac{11}{4} = 3$.
	} 
\end{ex}

% nguồn: Ok_Toan11_TracNghiem_Lan2
\begin{ex}%[1D2H3-6]%[ma:9fbcef8e]
	Một cấp số nhân $(u_n)$ có $n$  số hạng, số hạng đầu $u_1 = 7$, công bội $q = 2$. Số hạng thứ $n$ bằng $1\,792$. Tính tổng $n$ số hạng đầu tiên của cấp số nhân $(u_n)$.
	\choice 
		{$5\,377$} 
		{$5\,737$} 
		{\True  $3\,577$}
		{$3\,775$} 
	 \loigiai{
		Ta  có $u_n = u_1 \cdot q^{n-1} \Rightarrow 7 \cdot 2^{n-1} = 1\,792 \Leftrightarrow n = 9 \Rightarrow S_9 = 3\,577$.
	} 
\end{ex}

% nguồn: Ok_Toan11_TracNghiem_Lan2
\begin{ex}%[1D2H3-6]%[ma:08213391]
	Cho dãy $\left(u_n\right)$ với $u_n =  \left(\dfrac{1}{2}\right)^n + 1,\ \forall n \in \mathbb{N}^*$. Tính $S_{2019} = u_1 + u_2 + u_3 + \ldots + u_{2019}$, ta được kết quả
	\choice 
		{\True  $2020 - \dfrac{1}{2^{2019}}$}
		 {$\dfrac{4039}{2}$}
		{$2019 + \dfrac{1}{2^{2019}}$} 
		 {$\dfrac{6057}{2}$}
	\loigiai{ 
		 Ta có:
		 \[
		S_{2019} = 2019 + \left(\dfrac{1}{2}\right)^1 + \left(\dfrac{1}{2}\right)^2 + \ldots + \left(\dfrac{1}{2}\right)^{2019} = 2019 + \dfrac{1}{2}\cdot\dfrac{1 - \left(\dfrac{1}{2}\right)^{2019}}{1 -  \dfrac{1}{2}} = 2020 - \dfrac{1}{2^{2019}}.
		 \]
	} 
\end{ex}

% nguồn: Ok_Toan11_TracNghiem_Lan2
\begin{ex}%[1D2H3-6]%[ma:dab278c3]
	Cho dãy số $(a_n)$ xác định bởi $a_1 = 2,\ a_{n+1} = -2a_n,\ n \ge 1,\ n \in \mathbb{N}$. Tính  tổng của $10$ số hạng đầu tiên của dãy số.
	 \choice
		 {$\dfrac{2050}{3}$}
		{$2046$} 
		{\True  $-682$}
		{$-2046$} 
	\loigiai{ 
		Vì $\dfrac{a_{n+1}}{a_n} = -2$ suy  ra $(a_n)$ là một cấp số nhân với
		$\heva{ 
		&a_1  = 2\\
		 &q = -2.
		 }$\\
		 Suy ra $S_{10} = \dfrac{a_1\left(1 - q^{10}\right)}{1 - q} = \dfrac{2(1 - (-2)^{10})}{1 + 2} = -682$.
	 }
\end{ex}

% nguồn: Ok_Toan11_TracNghiem_Lan2
\begin{ex}%[1D2H3-6]%[ma:56973c9d]
	Cho dãy số $(u_n)$,  $n\in \mathbb{N}^{*}$ thỏa mãn điều kiện
	$\heva{&u_1=3\\ 
	&u_{n+1}=-\dfrac{u_n}{5}}$. Gọi $S=u_1+u_2+u_3+\cdots+u_n$ là tổng $n$ số hạng đầu tiên của dãy  số đã cho. Khi đó $\lim S_n$ bằng
	\choice 
		 {$\dfrac{1}{2}$}
		 {$\dfrac{3}{5}$}
		{$0$} 
		{\True  $\dfrac{5}{2}$}
	\loigiai{Ta có $\dfrac{u_{n+1}}{u_n}=\dfrac{-\dfrac{u_n}{5}}{u_n}=-\dfrac{1}{5}$ do đó dãy $(u_n), n\in \mathbb{N}^{*}$ là một cấp số nhân lùi vô hạn  có $u_1=3$
		,  $d=-\dfrac{1}{5}$.\\
		Suy  ra $\lim S_n=\dfrac{u_1}{1-q}=\dfrac{3}{1+\dfrac{1}{5}}=\dfrac{5}{2}$.
	 }
\end{ex}

% nguồn: Ok_Toan11_TracNghiem_Lan2
\begin{ex}%[1D2H3-6]%[ma:36d7c17d]
	Số thập phân vô hạn tuần hoàn $3{,}15555\ldots=3{,}1(5)$ viết dưới dạng hữu tỉ  là
	 \choice
		 {$\dfrac{63}{20}$}
		 {\True$\dfrac{142}{45}$}
		{$\dfrac{1}{18}$} 
		{$\dfrac{7}{2}$} 
	\loigiai{\[3{,}15555\ldots=3{,}1(5)=3{,}1+5\left(\dfrac{1}{10^2}+\dfrac{1}{10^3}+\cdots\right)=3{,}1+5\cdot\dfrac{\dfrac{1}{10^2}}{1-\dfrac{1}{10}}=\dfrac{142}{45}. 
		 \]
	} 
\end{ex}

% nguồn: Ok_Toan11_TracNghiem_Lan2
\begin{ex}%[1D2H3-6]%[ma:45d9d61c]
	Cho cấp số nhân $\left(u_n\right)$ có  $\heva{&u_4+u_6=-540\\&u_3+u_5=180}$. Tính $S_{21}$.
	 \choice
		{\True  $S_{21}=\dfrac{1}{2}\left(3^{21}+1\right)$}
		{$S_{21}=3^{21}-1$} 
		 {$S_{21}=1-3^{21}$}
		{$S_{21}=-\dfrac{1}{2}\left(3^{21}+1\right)$} 
	\loigiai{ 
		Ta có $u_4+u_6=-540$ $\Leftrightarrow\left(u_3+u_5\right)q=-540$.\\ 
		Kết hợp với phương trình thứ hai  trong hệ, ta tìm được $q=-3$.\\
		Lại có $u_3+u_5=180\Leftrightarrow u_1\left(q^2+q^4\right)=180$.\\ 
		 Vì $q=-3$ nên $u_1=2$.\\
		 Suy ra $S_{21}=\dfrac{u_1\left(1-q^{21}\right)}{1-q}=\dfrac{1}{2}\left(3^{21}+1\right)$.\\
	} 
\end{ex}

% nguồn: Ok_Toan11_TracNghiem_Lan2
\begin{ex}%[1D2H3-6]%[ma:b33a99f2]
	Cho cấp số nhân $\left(u_n\right)$ thỏa mãn $\heva{&u_1+u_2+u_3=13\\&u_4-u_1=26}$. Tổng $8$ số  hạng đầu của cấp số nhân $\left(u_n\right)$ là
	\choice 
		{$S_8=1\,093$} 
		{$S_8=3\,820$} 
		{$S_8=9\,841$} 
		{\True $S_8=3\,280$} 
	\loigiai{ 
		 Ta có
		 \allowdisplaybreaks
		\begin{eqnarray*}
			 &&\heva{&u_1+u_2+u_3=13\\&u_4-u_1=26}\\
			&\Leftrightarrow&\heva{&u_1+u_1\cdot q+u_1\cdot q^2=13\\& u_1\cdot q^3-u_1=26}\\ 
			&\Leftrightarrow&\heva{&u_1\left(1+q+q^2\right)=13\\&u_1\cdot  \left(q-1\right)\left(1+q+q^2\right)=26}\\
			&\Leftrightarrow&\heva{&u_1\left(1+q+q^2\right)=13\\&q=3}\\ 
			 &\Leftrightarrow&\heva{&u_1=1\\&q=3.}
		\end{eqnarray*}
		 Vậy tổng $S_8=\dfrac{u_1\left(1-q^8\right)}{1-q}=\dfrac{1\left(1-3^8\right)}{1-3}=3\,280$.
	 }
\end{ex}

% nguồn: Ok_Toan11_TracNghiem_Lan2
\begin{ex}%[1D2H3-6]%[ma:23246a19]
	Cho cấp số nhân $(u_n)$ có $\heva{&u_4 + u_6 = -540\\ &u_3 + u_5 = 180}$. Tính $S_{21}$. 
	 \choice
		 {\True $S_{21} = \dfrac{1}{2} \left( 3^{21} + 1 \right)$}
		{$S_{21} = 3^{21} -  1$}
		 {$S_{21} = 1 - 3^{21}$}
		{$S_{21} = -\dfrac{1}{2} \left( 3^{21} + 1  \right)$}
	 \loigiai{
		Ta có $u_4 + u_6  = -540 \Leftrightarrow (u_3 + u_5)q = -540$.\\
		Kết hợp với phương  trình thứ hai trong hệ, ta tìm được $q = -3$.\\
		Lại có $u_3 + u_5 = 180 \Leftrightarrow u_1(q^2  + q^4) = 180$.\\
		 Vì $q = -3$ nên $u_1 = 2$.\\
		Suy ra $S_{21} = \dfrac{u_1(1 - q^{21})}{1 - q}  = \dfrac{1}{2}(3^{21} + 1)$.
	 }
\end{ex}

% nguồn: Ok_Toan11_TracNghiem_Lan2
\begin{ex}%[1D2H3-6]%[ma:20ec0980]
	Cho cấp số nhân $(u_n)$ có $u_1 = 8$ và biểu thức $4u_3 + 2u_2  - 15u_1$ đạt giá trị nhỏ nhất. Tính $S_{10}$.
	\choice 
		 {$S_{10} = \dfrac{2\left(4^{11} + 1\right)}{5 \cdot 4^9}$}
		{\True $S_{10} = \dfrac{2\left(4^{10} - 1\right)}{5  \cdot 4^8}$}
		{$S_{10}  = \dfrac{2^{10} - 1}{3 \cdot 2^6}$}
		{$S_{10} = \dfrac{2^{11} -  1}{3 \cdot 2^7}$}
	 \loigiai{
		Gọi  $q$ là công bội của cấp số nhân. Khi đó $4u_3 + 2u_2 - 15u_1 = 2(4q + 1)^2 - 122 \geq -122,\ \forall q$.\\
		Dấu bằng xảy ra  khi $4q + 1 = 0 \Leftrightarrow q = -\dfrac{1}{4}$.\\
		 Suy ra
		\[ 
		S_{10} = u_1 \cdot \dfrac{1 -  q^{10}}{1 - q} = 8 \cdot \dfrac{1 - \left(-\dfrac{1}{4}\right)^{10}}{1 + \dfrac{1}{4}} = \dfrac{2\left(4^{10} - 1\right)}{5 \cdot 4^8}.
		\] 
	} 
\end{ex}

% nguồn: Ok_Toan11_TracNghiem_Lan2
\begin{ex}%[1D2H3-6]%[ma:d3c0a239]
	Một hình gồm các khối cầu xếp chồng lên nhau tạo thành một cột thẳng đứng. Biết rằng mỗi khỗi cầu có bán  kính gấp đôi bán kính của khối cầu nằm ngay trên nó và bán kính khối cầu dưới cùng là $50$ cm. Hỏi mệnh đề nào sau đây đúng.
	\choice 
		{Chiều cao của mô  hình không quá $1{,}5$ mét}
		{Chiều cao của mô  hình tối đa là $2$ mét}
		{\True Chiều cao của mô  hình dưới $2$ mét}
		{Mô hình  có thể đat được chiều cao tùy ý}
	\loigiai{ 
		Giả  sử có $n$ khối cầu.\\
		Gọi bán kính của  khối cầu thứ nhất (khối cầu dưới cùng) là $R_1=\dfrac{1}{2}$ mét.\\
		Gọi bán kính của khối cầu thứ $2$ (ngay trên khối  cầu thứ nhất) là $R_2$.\\
		Gọi bán kính của khối cầu thứ $n$  (ngay trên khối cầu thứ $n-1$) là $R_n$ ($n \in \mathbb{N}$, $n\geq 2$).\\
		Khi đó dãy các số $R_1$, $R_2$, $\ldots$,  $R_n$, $\ldots$ là một cấp số nhân với công bội $q=\dfrac{1}{2}$.\\
		Chiều cao của mô hình là 
		\begin{eqnarray*}
			h&=&2\left(R_1+R_2+\cdots+R_n\right)\\ 
			&=&2\cdot\dfrac{1}{2} \dfrac{1-\left(\dfrac{1}{2}\right)^n}{1-\dfrac{1}{2}}\\ 
			&=&2\left(1-\left(\dfrac{1}{2}\right)^n\right). 
		\end{eqnarray*}
		Suy ra với $h=2-\left(\dfrac{1}{2}\right)^{n-1} <  2$, $\forall n \in \mathbb{N}^{*}$.\\
		Do đó chiều cao của mô hình dưới $2$  mét.
	 }
\end{ex}

% nguồn: Ok_Toan11_TracNghiem_Lan2
\begin{ex}%[1D2H3-6]%[ma:20aca0dc]
	 Tổng $S = \dfrac{1}{3} + \dfrac{1}{3^2} + \cdots + \dfrac{1}{3^n} + \cdots$ có giá trị là
	\choice 
		 {$\dfrac{1}{9}$}
		{$\dfrac{1}{4}$} 
		 {$\dfrac{1}{3}$}
		 {\True $\dfrac{1}{2}$}
	\loigiai{ 
		Ta có $S = \dfrac{1}{3} + \dfrac{1}{3^2} + \cdots + \dfrac{1}{3^n} + \cdots$ là tổng của cấp số nhân lùi vô  hạn $(u_n)$ với $u_n = \dfrac{1}{3^n}$ có số hạng đầu $u_1 = \dfrac{1}{3}$, công sai $q = \dfrac{1}{3}$.\\
		Do đó  $S = \dfrac{u_1}{1 - q} = \dfrac{\dfrac{1}{3}}{1 - \dfrac{1}{3}} = \dfrac{1}{2}$.
	} 
\end{ex}

% nguồn: Ok_Toan11_TracNghiem_Lan2
\begin{ex}%[1D2H3-6]%[ma:2b662b42]
	Tổng vô hạn sau  đây $S=2+\dfrac{2}{3}+\dfrac{2}{3^2}+\cdots+\dfrac{2}{3^n}+\cdots$ có giá trị bằng
	 \choice
		 {$\dfrac{8}{3}$}
		{\True  $3$}
		{$4$} 
		{$2$} 
	\loigiai{Ta có $2;\dfrac{2}{3};\dfrac{2}{3^2};\ldots;\dfrac{2}{3^n};\ldots$ là một cấp số nhân lùi vô hạn với công bội $q=\dfrac{1}{3}  < 1$
		 \[S=2+\dfrac{2}{3}+\dfrac{2}{3^2}+\cdots+\dfrac{2}{3^n}+\cdots=2.\dfrac{1}{1-\dfrac{1}{3}}=3.
		 \]
	} 
\end{ex}

% nguồn: Ok_Toan11_TracNghiem_Lan2
\begin{ex}%[1D2H3-6]%[ma:8312c098]
	Cho cấp  số nhân $(u_n)$ có $S_2 = 4$ và $S_3 = 13$. Tìm $S_5$.
	\choice 
		{\True $S_5 = 121$ hoặc  $S_5 = \dfrac{181}{16}$}
		{$S_5  = 121$ hoặc $S_5 = \dfrac{35}{16}$}
		{$S_5 = 114$ hoặc $S_5 =  \dfrac{185}{16}$}
		{$S_5 = 141$ hoặc $S_5  = \dfrac{183}{16}$}
	 \loigiai{
		Ta có $u_3  = S_3 - S_2 = 9 \Rightarrow u_1 q^2 = 9 \Rightarrow u_1 = \dfrac{9}{q^2}$.\\
		Vì $S_2 = 4$ nên $u_1 + u_1 q = 4$. Do đó $\dfrac{9}{q^2} + \dfrac{9}{q} = 4 \Leftrightarrow  4q^2 - 9q - 9 = 0 \Leftrightarrow q = 3$ hoặc $q = -\dfrac{3}{4}$.\\
		Với $q = 3$ thì $u_1 = 1$, $u_6 =  u_1 q^5 = 243$. Suy ra $S_5 = \dfrac{u_1 - u_6}{1 - q} = \dfrac{1 - 243}{1 - 3} = 121$.\\
		Với $q = -\dfrac{3}{4}$ thì $u_1 = 16$, $u_6 = -\dfrac{243}{64}$. Suy ra $S_5 = \dfrac{u_1 - u_6}{1 - q}  = \dfrac{181}{16}$.
	 }
\end{ex}

% nguồn: Ok_Toan11_TracNghiem_Lan2
\begin{ex}%[1D2H3-6]%[ma:f61ee2df]
	Cho dãy số $\left(a_n\right)$ xác định bởi $a_1=2$, $a_{n+1}=-2a_n$,  $n\ge 1$, $n\in\mathbb{N}$. Tính tổng của $10$ số hạng đầu tiên của dãy số.
	 \choice
		{$\dfrac{2\,050}{3}$} 
		{$2\,046$} 
		{\True $-682$} 
		 {$-2\,046$}
	\loigiai{ 
		Vì $\dfrac{a_{n+1}}{a_n}=-2$ suy ra $\left(a_n\right)$  là một cấp số nhân với $\heva{&a_1=2\\&q=-2.}$\\
		Suy ra  $S_{10}=\dfrac{a_1\left(1-q^{10}\right)}{1-q}=-682$.
	 }
\end{ex}

% nguồn: Ok_Toan11_TracNghiem_Lan2
\begin{ex}%[1D2H3-6]%[ma:6f0bd22d]
	 Tổng $1+\dfrac{1}{2}+\dfrac{1}{4}+\dfrac{1}{2^n}+\cdots$ bằng
	\choice 
		{$\dfrac{1}{2}$} 
		 {\True$2$}
		 {$1$}
		{$+\infty$} 
	\loigiai{Ta có $1+\dfrac{1}{2}+\dfrac{1}{4}+\dfrac{1}{2^n}+\cdots$ là tổng của một cấp số nhân lùi vô hạn với $u_1=1$, $q=\dfrac{1}{2}$.\\ 
		Áp dụng công thức được $S=\dfrac{u_1}{1-q}$  kết quả $1+\dfrac{1}{2}+\dfrac{1}{4}+\dfrac{1}{2^n}+\cdots=2$.
	 }
\end{ex}

% nguồn: Ok_Toan11_TracNghiem_Lan2
\begin{ex}%[1D2H3-6]%[ma:87e53aa2]
	 Tính tổng $S$ của cấp số nhân lùi vô hạn có số hạng đầu $u_1=1$ và công bội $q=-\dfrac{1}{2}$.
	 \choice
		 {$S=2$}
		{\True $S=\dfrac{3}{2}$} 
		{$S=1$} 
		{$S=\dfrac{2}{3}$} 
	 \loigiai{\[S=\dfrac{u_1}{1-q}=\dfrac{1}{1+\dfrac{1}{2}}=\dfrac{2}{3}.
		\] 
	 }
\end{ex}

% nguồn: Ok_Toan11_TracNghiem_Lan2
\begin{ex}%[1D2H3-6]%[ma:b8f445cb]
	Số thập phân vô hạn tuần hoàn $0{,}121212\ldots$ được biểu  diễn bởi phân số
	 \choice
		{$\dfrac{3}{25}$} 
		 {\True $\dfrac{12}{99}$}
		{$\dfrac{1}{11}$} 
		 {$\dfrac{3}{22}$}
	\loigiai{ 
		 Ta có
		\begin{align*}
			0{,}121212\ldots 
			&= \dfrac{12}{10^2} + \dfrac{12}{10^4} + \dfrac{12}{10^6}  + \ldots + \dfrac{12}{10^{2n}} + \ldots\\
			&= 12\left( \dfrac{1}{10^2} + \dfrac{1}{10^4} + \ldots + \dfrac{1}{10^{2n}} + \ldots  \right) \\
			&=12 \left( \dfrac{\tfrac{1}{100} }{1  - \tfrac{1}{100}} \right) \\
			 &= \dfrac{4}{33} = \dfrac{12}{99}.
		\end{align*}
	 }
\end{ex}

% nguồn: Ok_Toan11_TracNghiem_Lan2
\begin{ex}%[1D2H3-6]%[ma:6737a513]
	Tổng các số nhân  lùi vô hạn $- \dfrac{1}{2},\ -\dfrac{1}{4},\ -\dfrac{1}{8},\ldots,\ \dfrac{-1}{2^n},\ldots$ là
	 \choice
		{\True $-1$} 
		 {$\dfrac{1}{2}$}
		 {$-\dfrac{1}{4}$}
		 {$-\dfrac{1}{3}$}
	 \loigiai{
		Cấp số nhân có $u_1 = -\dfrac{1}{2}$, công bội $q = \dfrac{1}{2}$ nên tổng của cấp  số nhân lùi vô hạn là $S_n = \dfrac{u_1}{1 - q} = -1$.
	} 
\end{ex}

% nguồn: Ok_TracNghiemLop11_Lan1
\begin{ex}%[1D2H3-6]%[ma:b97121b9]
	Cho năm số $a,\, b,\, c,\,  d,\, e$ tạo thành một cấp số nhân theo thứ tự đó và các số đều khác $0$, biết $\dfrac{1}{a}+\dfrac{1}{b}+\dfrac{1}{c}+\dfrac{1}{d}+\dfrac{1}{e}=10$ và tổng của chúng bằng $40$. Tính giá trị $|S|$ với $S=a b c d e$.
	 \choice
		 {$|S|=52$}
		 {$|S|=42$}
		{$|S|=62$} 
		 {\True $|S|=32$}
	 \loigiai{
		Gọi $q(q \neq 0)$ là công bội của cấp số  nhân $a,\, b,\, c,\, d,\, e$. Khi đó $\dfrac{1}{a},\, \dfrac{1}{b},\, \dfrac{1}{c},\, \dfrac{1}{d}$,\, $\dfrac{1}{e}$ là cấp số nhân có công bội $\dfrac{1}{q}$.\\
		Theo đề bài ta  có
		$$ 
		 \heva{
		&a + b + c + d + e = 4  0\\
		&\dfrac { 1 }{ a } + \dfrac { 1 }{ b } + \dfrac {  1 }{ c } + \dfrac { 1 }{ d } + \dfrac { 1 }{ e } = 1 0
		} 
		\Leftrightarrow \heva{ 
		&a \cdot \dfrac { 1 - q ^ { 5 } }{  1 - q } = 4 0 \\
		&\dfrac { 1 - ( \dfrac { 1 }{ q }) ^ { 5 } }{ a }  \cdot \dfrac { 1 - \dfrac { 1 }{ q } }{ 1 } = 1 0
		} 
		\Leftrightarrow \heva{ 
		&a \cdot  \dfrac{1-q^5}{1-q}=40 \\
		&\dfrac{1}{a} \cdot  \dfrac{q^5-1}{q^4(q-1)}=10
		} 
		\Leftrightarrow a^2  q^4=4.
		$$ 
		Ta có $S=a b c d e=a \cdot a q \cdot  a q^2 \cdot a q^3 \cdot a q^4=a^5 q^{10}$.\\
		Nên $S^2=\left(a^5 q^{10}\right)^2=\left(a^2  q^4\right)^5=4^5$.\\
		 Suy ra $|S|=\sqrt{4^5}=32$.\\
	 }
\end{ex}

% nguồn: Ok_TracNghiemLop11_Lan1
\begin{ex}%[1D2H3-6]%[ma:eed78130]
	Tổng vô hạn $S=1+\dfrac{1}{2}+\dfrac{1}{2^2}+\ldots+\dfrac{1}{2^n}+\ldots$ bằng 
	 \choice
		 {$4$}
		 {$2^n-1$}
		 {$1$}
		{\True  $2$}
	 \loigiai{
		Đây  là tổng của một cấp số nhân lùi vô hạn, với $u_1=1$; $q=\dfrac{1}{2}$.\\
		Khi đó  $S=\dfrac{u_1}{1-q}$ $=\dfrac{1}{1-\dfrac{1}{2}}$ $=2$.
	} 
\end{ex}

% nguồn: Ok_Toan11_TracNghiem_Lan2
\begin{ex}%[1D2V3-6]%[ma:05c0f698]
	Một hình vuông $ABCD$ có cạnh $AB=a$, diện tích $S_1$. Nối $4$ trung  điểm $A_1$, $B_1$, $C_1$, $D_1$ theo thứ tự của $4$ cạnh $AB$, $BC$, $CD$, $DA$ ta được hình vuông thứ hai là $A_1B_1C_1D_1$ có diện tích $S_2$. Tiếp tục như thế ta được hình vuông thứ ba $A_2B_2C_2D_2$ có diện tích $S_3$ và cứ tiếp tục như thế, ta được diện tích $S_4$, $S_5$, $\ldots$ Tính $S=S_1+S_2+S_3+\cdots+S_{100}$.
	 \choice
		{$S=\dfrac{2^{100}-1}{2^{99}  a^2}$}
		{$S=\dfrac{a\left(2^{100}-1\right)}{2^{99}}$} 
		{\True $S=\dfrac{a^2\left(2^{100}-1\right)}{2^{99}}$} 
		{$S=\dfrac{a^2\left(2^{99}-1\right)}{2^{99}}$} 
	 \loigiai{
		Dễ thấy $S_1=a^2$; $S_2=\dfrac{a^2}{2}$; $S_3=\dfrac{a^2}{4}$;  $\ldots$; $S_{100}=\dfrac{a^2}{2^{99}}$.\\
		Như vậy $S_1$, $S_2$, $S_3$, $\ldots$, $S_{100}$ là cấp số nhân với công bội $q=\dfrac{1}{2}$.\\ 
		$S=S_1+S_2+\cdots+S_{100}=a^2\left(1+\dfrac{1}{2}+\dfrac{1}{2^2}+\cdots+\dfrac{1}{2^{99}}  \right)=\dfrac{a^2\left(2^{100}-1\right)}{2^{99}}$.
	} 
\end{ex}

% nguồn: Ok_Toan11_TracNghiem_Lan2
\begin{ex}%[1D2V3-6]%[ma:7f5b08d5]
	Cho dãy số $(U_n)$ xác định bởi $U_1=\dfrac{1}{3}$ và  $U_{n+1}=\dfrac{n+1}{3n}U_n$. Tổng $S=U_1+\dfrac{U_2}{2}+\dfrac{U_3}{3}+\dots+\dfrac{U_{10}}{10}$ bằng
	\choice 
		 {$\dfrac{3280}{6561}$}
		 {\True $\dfrac{29524}{59049}$}
		 {$\dfrac{25942}{59049}$}
		{$\dfrac{1}{243}$} 
	\loigiai{ 
	Theo  đề ta có $U_{n+1} = \dfrac{n+1}{3n}U_n \Leftrightarrow \dfrac{U_{n+1}}{n+1} = \dfrac{1}{3}\cdot\dfrac{U_n}{n}$ mà $U_1 = \dfrac{1}{3}$ hay $\dfrac{U_1}{1} = \dfrac{1}{3}$.\\
	Nên ta có $\dfrac{U_2}{2} = \dfrac{1}{3} \cdot\dfrac{U_1}{1} = \dfrac{1}{3}\cdot \left(\dfrac{1}{3} \right) = \left(\dfrac{1}{3} \right)^2$;$\dfrac{U_3}{3} = \dfrac{1}{3}\cdot  \dfrac{U_2}{2} = \dfrac{1}{3}\cdot \left(\dfrac{1}{3} \right)^2 = \left(\dfrac{1}{3} \right)^3$;$\ldots$;$\dfrac{U_{10}}{10} = \left(\dfrac{1}{3} \right)^{10}$.\\
	Hay dãy $\left(\dfrac{U_n}{n} \right)$ là một cấp số nhân có số hạng đầu $U_1 = \dfrac{1}{3}$, công bội  $q = \dfrac{1}{3}$.\\
	Khi đó $S = U_1 + \dfrac{U_2}{2} + \dfrac{U_3}{3} + \dots + \dfrac{U_{10}}{10} = \dfrac{3^{10} -  1}{2\cdot3^{10}} = \dfrac{59048}{2\cdot3^{10}} = \dfrac{29524}{59049}$.}
\end{ex}

% nguồn: Ok_Toan11_TracNghiem_Lan2
\begin{ex}%[1D2V3-6]%[ma:e1525c84]
	Cho cấp số nhân $(a_n)$ có $a_1=7$, $a_6=224$ và $S_k=3\,577$. Tính  giá trị của biểu thức $T=(k+1) a_k$.
	\choice 
		{\True  $T=17\,920$}
		{$T=8\,064$} 
		{$T=39\,424$} 
		{$T=86\,016$} 
	\loigiai{Ta  có $a_6=224 \Leftrightarrow a_1q^5=224 \Rightarrow q=2$.\\
		Do $S_k=\dfrac{a_1(1-q^k)  }{1-q}=7(2^k-1)$ nên $S_k=3\,577 \Leftrightarrow 7(2^k-1) =3\,577 \Leftrightarrow 2^k=2^9 \Leftrightarrow k=9$.\\
	Suy ra $T=10a_9=10a_1q^8=17\,920$.} 
\end{ex}

% nguồn: Ok_Toan11_TracNghiem_Lan2
\begin{ex}%[1D2V3-6]%[ma:e1100679]
	Giá trị của tổng $7+77+777+\cdots+77\ldots 7$ bằng (với $77\ldots 7$ có  $2\,018$ số $7$).
	 \choice
		 {$\dfrac{70}{9} \left(10^{2\,018}-1\right)+2\,018$}
		 {$\dfrac{7}{9} \left(\dfrac{10^{2\,018}-10}{9}-2\,018\right)$}
		{\True $\dfrac{7}{9} \left(\dfrac{10^{2\,019}-10}{9}-2\,018\right)$} 
		{$\dfrac{7}{9} \left(10^{2\,018}-1\right)$} 
	 \loigiai{
		Ta có 
		\begin{eqnarray*}
			7+77+777+\cdots+77\ldots 7 
			&=&\dfrac{7}{9}  \left(9+99+999+\cdots+99\ldots 9\right)\\
			 &=&\dfrac{7}{9} \left(10-1+10^2-1+10^3-1+\cdots+10^{2\,018}-1\right)\\
			 &=&\dfrac{7}{9} \left(10+10^2+10^3+\cdots+10^{2\,018}-2\,018\right).
		\end{eqnarray*}
		Mặt khác,ta có $10+10^2+10^3+\cdots+10^{2\,018}$  là tổng của một cấp số nhân với $u_1=10$ và công bội $q=10$\\
		Suy ra  $10+10^2+10^3+\cdots+10^{2\,018} =10\dfrac{10^{2\,018}-1}{9}=\dfrac{10^{2\,019}-10}{9}$.\\
		Do  đó $\dfrac{7}{9} \left(10+10^2+10^3+\cdots+10^{2\,018}-2\,018\right)=\dfrac{7}{9} \left(\dfrac{10^{2\,019}-10}{9}-2\,018\right)$.
	 }
\end{ex}

% nguồn: Ok_Toan11_TracNghiem_Lan2
\begin{ex}%[1D2V3-6]%[ma:97edc6cd]
	Cho dãy số $(u_n)$ là một cấp số nhân có số hạng  đầu $u_1=1$, công bội $q=2$. Tính tổng $T=\dfrac{1}{u_1-u_5}+\dfrac{1}{u_2-u_6}+\dfrac{1}{u_3-u_7}+\ldots+\dfrac{1}{u_{20}-u_{24}}$.
	 \choice
		 {$\dfrac{1-2^{19}}{15\cdot2^{18}}$}
		 {\True $\dfrac{1-2^{20}}{15\cdot2^{19}}$}
		{$\dfrac{2^{19}-1}{15\cdot2^{18}}$} 
		{$\dfrac{2^{20}-1}{15\cdot2^{19}}$} 
	\loigiai{\begin{align*}
		 T
		&=\dfrac{1}{u_1-u_5}+\dfrac{1}{u_2-u_6}+\dfrac{1}{u_3-u_7}+\ldots+\dfrac{1}{u_{20}-u_{24}}  \\
		&=\dfrac{1}{u_1(1-q^4)  }+\dfrac{1}{u_2(1-q^4) }+\dfrac{1}{u_3(1-q^4) }+\ldots+\dfrac{1}{u_{20}(1-q^4) }\\
		&=\dfrac{1}{1-q^4}\left(\dfrac{1}{u_1}+\dfrac{1}{u_2}+\dfrac{1}{u_3}+\ldots+\dfrac{1}{u_{20}}\right)  \\
		 &=\dfrac{1}{1-q^4}\left(\dfrac{1}{u_1}+\dfrac{1}{u_1q}+\dfrac{1}{u_1q^2}+\ldots+\dfrac{1}{u_1q^{19}}\right)
		 =\dfrac{1}{1-q^4}\cdot\dfrac{1}{u_1}\left(1+\dfrac{1}{q}+\dfrac{1}{q^2}+\ldots+\dfrac{1}{q^{19}}\right)\\
		&=\dfrac{1}{1-q^4}\cdot\dfrac{1}{u_1}\dfrac{1-\left(\dfrac{1}{q}\right)^{20}}{1-\dfrac{1}{q}}  =\dfrac{1}{1-q^4}\cdot\dfrac{1-(q)^{20}}{u_1(1-q) q^{19}}\\
		 &=\dfrac{1-2^{20}}{15\cdot2^{19}}.
	\end{align*}
 }
\end{ex}

% nguồn: Ok_Toan11_TracNghiem_Lan2
\begin{ex}%[1D2V3-6]%[ma:ea6c5edc]
	Tính tổng tất cả các số hạng của một cấp số nhân có số hạng đầu là $\dfrac{1}{2}$, số hạng  thứ tư là $32$ và số hạng cuối là $2\,048$.
	\choice 
		{$\dfrac{1\,365}{2}$} 
		 {$\dfrac{5\,416}{2}$}
		{\True $\dfrac{5\,461}{2}$} 
		 {$\dfrac{21\,845}{2}$}
	 \loigiai{
		 Theo bài ra ta có $u_1=\dfrac{1}{2}$, $u_4=32$ và $u_n=2\,048$.\\
		$u_4=u_1\cdot q^3\Rightarrow 32=\dfrac{1}{2}\cdot  q^3\Rightarrow q=4$.\\
		$u_n=2\,048\Rightarrow  u_1\cdot q^{n-1}=2\,048\Rightarrow 4^{n-1}=4^6\Rightarrow n=7$.\\
		Khi đó tổng của cấp số nhân này là $S_7=\dfrac{u_1 \left(1-q^7\right)}{1-q}=\dfrac{\dfrac{1}{2} \left(1-4^7\right)}{1-4}=\dfrac{5\,461}{2}$. 
	} 
\end{ex}

% nguồn: Ok_Toan11_TracNghiem_Lan2
\begin{ex}%[1D2V3-6]%[ma:8353a469]
	 Cho cấp số nhân $(u_n)$ có $u_1 = 2$, công bội dương và biểu thức $u_4 + \dfrac{1\,024}{u_7}$ đạt giá trị nhỏ nhất.
	Tính $S = u_{11} +  u_{12} + \ldots + u_{20}$.
	 \choice
		 {$S = 2\,046$}
		{\True $S = 2\,095\,104$} 
		{$S =  2\,097\,150$}
		{$S = 1\,047\,552$} 
	 \loigiai{
		Gọi $q$ là công bội của cấp số nhân, $q > 0$. Ta có  $u_4 + \dfrac{1\,024}{u_7} = 2q^3 + \dfrac{512}{q^6}$.\\
		Áp dụng bất đẳng thức Cauchy,  ta có
		 \[
		2q^3 + \dfrac{512}{q^6} = q^3 + q^3 + \dfrac{512}{q^6} \geq 3 \sqrt[3]{q^3 \cdot  q^3 \cdot \dfrac{512}{q^6}} = 24.
		\] 
		Suy ra $u_4  + \dfrac{1\,024}{u_7}$ đạt giá trị nhỏ nhất bằng $24$ khi $q^3 = \dfrac{512}{q^6} \Leftrightarrow q = 2$.\\
		Ta có: 
		 \[
		S_{10} = \dfrac{u_1(1 - q^{10})}{1 - q} = \dfrac{2(1 - 2^{10})}{1 - 2} = 2^{11} - 2, \ S_{20} = \dfrac{u_1(1 -  q^{20})}{1 - q} = \dfrac{2(1 - 2^{20})}{1 - 2} = 2^{21} - 2.
		 \]
		Do đó $S = S_{20}  - S_{10} = 2\,095\,104$.
	} 
\end{ex}

% nguồn: Ok_Toan11_TracNghiem_Lan2
\begin{ex}%[1D2V3-6]%[ma:3a4d4101]
	Cho dãy số $(u_n)$ thỏa mãn $\heva{& u_1=1 \\& u_n=2u_{n-1}+1;n \geq 2}$. Tổng $S=u_1+u_2+\dots+u_{20}$ bằng 
	 \choice
		 {$2^{20}-20$}
		 {\True $2^{21}-22$}
		 {$2^{20}$}
		 {$2^{21}-20$}
	\loigiai{Ta có $u_n = 2u_{n-1} + 1  \Leftrightarrow u_n + 1 = 2(u_{n-1} + 1)$.\\
		 Đặt $v_n = u_n + 1$, ta có $v_n = 2v_{n-1}$ trong đó $v_1 = 2$.\\
		Vậy $(v_n) $ là cấp số nhân có số hạng đầu $v_1  = 2$ và công bội bằng $2$, nên số hạng tổng quát\\
		$v_n = 2^n \Rightarrow u_n = v_n - 1 =  2^n - 1$.\\
		Vậy  nên
		\begin{align*}
			S &= u_1  + u_2 + \dots + u_{20} =(2^1 - 1) +(2^2 - 1) + \dots +(2^{20} - 1) \\&=(2^1 + 2^2 + \dots + 2^{20}) - 20 \\&= 2\cdot(2^{20} - 1) - 20 \\&= 2^{21} - 22.
		\end{align*}
	 }
\end{ex}

% nguồn: Ok_Toan11_TracNghiem_Lan2
\begin{ex}%[1D2V3-6]%[ma:665f7ffb]
	Ba số phân biệt có tổng là $217$ có thể coi là các số hạng liên tiếp của một cấp số nhân, cũng có thể coi là số hạng thứ $2$, thứ $9$, thứ $44$ của một cấp số cộng. Hỏi phải lấy bao nhiêu số hạng đầu của cấp số cộng này để  tổng của chúng bằng $820$?
	 \choice
		 {\True $20$}
		 {$42$}
		 {$21$}
		{$17$} 
	\loigiai{ 
		Gọi ba số đó là $x$, $y$,  $z$.\\
		Do ba số là các số hạng thứ $2$, thứ $9$  và thứ $44$ của một cấp số cộng nên ta có $x$; $y=x+7d$; $z=x+42d$.\\
		Theo  giả thiết, ta có $x+y+z=x+x+7d+x+42d=3x+49d=217$.\\
		Mặt khác, do  $x$, $y$, $z$ là các số hạng liên tiếp của một cấp số nhân nên
		\[y^2=xz\Leftrightarrow  (x+7d)^2=x(x+42d)\Leftrightarrow d(-4x+7d)=0\Leftrightarrow \hoac{&d=0\\&-4x+7d=0.}\]
		Với $d=0$, ta có $x=y=z=\dfrac{217}{3}$. Suy ra  $n=820:\dfrac{217}{3}=\dfrac{2\,460}{217} \notin \mathbb{N}$.\\
		Với $-4x+7d=0$, ta có $\heva{&-4x+7d=0\\&3x+49d=217}\Leftrightarrow \hoac{&x=7\\&d=4}$.  Suy ra $u_1=7-4=3$.\\
		Do đó, $S_n=820\Leftrightarrow \dfrac{\left[2u_1+(n-1)d\right]n}{2}=820\Leftrightarrow  \dfrac{\left[2\cdot 3+4(n-1)\right]n}{2}=820\Leftrightarrow \hoac{&n=20\\&n=-\dfrac{41}{2}.}$\\
		Vậy  $n=20$.
	 }
\end{ex}

% nguồn: Ok_Toan11_TracNghiem_Lan2
\begin{ex}%[1D2V3-6]%[ma:a4bf845f]
	 Số thập phân vô hạn tuần hoàn $0{,}121\,212\ldots$ được biểu diễn bởi phân số
	\choice 
		 {$\dfrac{3}{25}$}
		{\True  $\dfrac{12}{99}$}
		{$\dfrac{1}{11}$} 
		{$\dfrac{3}{22}$} 
	 \loigiai{
		Ta có 
		 \allowdisplaybreaks
		\begin{eqnarray*}
			 0{,}121\,212\ldots&=&\dfrac{12}{10^2}+\dfrac{12}{10^4}+\dfrac{12}{10^6}+\cdots +\dfrac{12}{10^{2n}}+\cdots\\
			&=&12\left(\dfrac{1}{10^2}+\dfrac{1}{10^4}+\cdots  +\dfrac{1}{10^{2n}}+\cdots \right)\\
			&=&12\left(\dfrac{\dfrac{1}{100}}{1-\dfrac{1}{100}}\right)\\ 
			 &=&\dfrac{4}{33}\\
			 &=&\dfrac{12}{99}.
		\end{eqnarray*}
	 }
\end{ex}

% nguồn: Ok_Toan11_TracNghiem_Lan2
\begin{ex}%[1D2V3-6]%[ma:9d819f94]
	Biết rằng $S=1+2\cdot 3+3\cdot 3^2+\cdots +11\cdot 3^{10}=a+\dfrac{21\cdot  3^b}{4}$. Tính $P=a+\dfrac{b}{4}$.
	 \choice
		 {$P=1$}
		{$P=2$} 
		{\True  $P=3$}
		{$P=4$} 
	\loigiai{ 
		Từ  giả thiết suy ra $3S=3+2\cdot 3^2+3\cdot 3^3+\cdots +11\cdot 3^{11}$.\\
		Do đó $-2S=S-3S=1+3+3^2+\cdots +3^{10}-10\cdot  3^{11}=\dfrac{1-3^{11}}{1-3}-11\cdot 3^{11}=-\dfrac{1}{2}-\dfrac{21\cdot 3^{11}}{2}\Rightarrow S=\dfrac{1}{4}+\dfrac{21}{4}\cdot {3^{11}}$.\\
		Vì $S=\dfrac{1}{4}+\dfrac{21\cdot 3^{11}}{4}=a+\dfrac{21\cdot 3^b}{4}\Rightarrow a=\dfrac{1}{4}$, $b=11$. 
		Suy ra $P=\dfrac{1}{4}+\dfrac{11}{4}=3$. 
	 }
\end{ex}

% nguồn: Ok_Toan11_TracNghiem_Lan2
\begin{ex}%[1D2V3-6]%[ma:a95d4e36]
	Viết thêm bốn số vào giữa hai số $160$ và $5$ để được  một cấp số nhân. Tổng các số hạng của cấp số nhân đó là
	 \choice
		{$215$} 
		{\True  $315$}
		{$415$} 
		{$515$} 
	\loigiai{ 
		Từ giả  thiết ta có $\heva{&u_1=160\\&u_6=5}\Rightarrow q=\sqrt[5]{\dfrac{u_6}{u_1}}=\dfrac{1}{2}$.\\
		Suy ra tổng các số hạng của  cấp số nhân đó là $S=\dfrac{u_1\left(1-q^6\right)}{1-q}=\dfrac{160\left(1-\left(\dfrac{1}{2}\right)^6\right)}{\dfrac{1}{2}}=315$.
	} 
\end{ex}

% nguồn: Ok_Toan11_TracNghiem_Lan2
\begin{ex}%[1D2V3-6]%[ma:81234355]
	 Giá trị của tổng $7 + 77 + 777 + \ldots + \underbrace{77\ldots7}_{2\,018\text{ số }7} $ bằng
	\choice 
		{$\dfrac{70}{9}\left(10^{2\,018} -  1\right) + 2\,018$}
		{$\dfrac{7}{9}\left(\dfrac{10^{2\,018} -  10}{9} - 2\,018\right)$}
		{\True $\dfrac{7}{9}\left(\dfrac{10^{2019}  - 10}{9} - 2\,018\right)$}
		{$\dfrac{7}{9}\left(10^{2\,018}  - 1\right)$}
	\loigiai{ 
		Ta  có
		\begin{align*}
			7 + 77 + 777  + \ldots + \underbrace{77\ldots7}_{2\,018\text{ số }7}
			&= \dfrac{7}{9}\left(9 + 99 + 999 + \ldots + \underbrace{99\ldots9}_{2\,018\text{  số }9}\right) \\
			&= \dfrac{7}{9}\left(10 -  1 + 10^2 - 1 + 10^3 - 1 + \ldots + 10^{2\,018} - 1\right)\\
			&= \dfrac{7}{9}\left(10  + 10^2 + 10^3 + \ldots + 10^{2\,018} - 2\,018\right).
		\end{align*}
		Mặt khác, ta có $10 + 10^2 + 10^3 + \ldots + 10^{2\,018}$ là tổng  của một cấp số nhân với $u_1 = 10$ và công bội $q = 10$.\\
		Do đó $10 + 10^2 + 10^3 + \ldots  + 10^{2\,018} = 10\cdot \dfrac{10^{2\,018} - 1}{9} = \dfrac{10^{2019} - 10}{9}$.\\
		Suy ra $7 + 77 + 777 + \ldots  + \underbrace{77\ldots7}_{2\,018\text{ số }7} = \dfrac{7}{9}\left(\dfrac{10^{2019} - 10}{9} - 2\,018\right)$.
	} 
\end{ex}

% nguồn: Ok_Toan11_TracNghiem_Lan2
\begin{ex}%[1D2V3-6]%[ma:2e155090]
	Cho dãy số $\left(u_n \right)$ là một cấp số nhân có số hạng đầu $u_1=1$,  công bội $q=2$. Tính tổng $T=\dfrac{1}{u_1-u_5}+\dfrac{1}{u_2-u_6}+\dfrac{1}{u_3-u_7}+\cdots+\dfrac{1}{u_{20}-u_{24}}\cdot$
	 \choice
		{$\dfrac{1-2^{19}}{15\cdot  2^{18}}$}
		{\True  $\dfrac{1-2^{20}}{15\cdot 2^{19}}$}
		 {$\dfrac{2^{19}-1}{15\cdot 2^{18}}$}
		{$\dfrac{2^{20}-1}{15\cdot 2^{19}}$} 
	\loigiai{ 
		\begin{eqnarray*} T
			 &=&\dfrac{1}{u_1-u_5}+\dfrac{1}{u_2-u_6}+\dfrac{1}{u_3-u_7}+\cdots+\dfrac{1}{u_{20}-u_{24}} \\
			&=&\dfrac{1}{u_1 \left(1-q^4\right)}+\dfrac{1}{u_2 \left(1-q^4\right)}+\dfrac{1}{u_3 \left(1-q^4\right)}+\cdots+\dfrac{1}{u_{20} \left(1-q^4\right)}\\ &=&\dfrac{1}{1-q^4} \left(\dfrac{1}{u_1}+\dfrac{1}{u_2}+\dfrac{1}{u_3}+\cdots+\dfrac{1}{u_{20}} \right) \\ 
			&=&\dfrac{1}{1-q^4} \left(\dfrac{1}{u_1}+\dfrac{1}{u_1 q}+\dfrac{1}{u_1 q^2}+\cdots+\dfrac{1}{u_1 q^{19}} \right) \\ &=&\dfrac{1}{1-q^4}\cdot\dfrac{1}{u_1} \left(1+\dfrac{1}{q}+\dfrac{1}{q^2}+\cdots+\dfrac{1}{q^{19}} \right) \\ &=&\dfrac{1}{1-q^4}\cdot\dfrac{1}{u_1}\cdot\dfrac{\left(\dfrac{1}{q}  \right)^{20}-1}{\dfrac{1}{q}-1}\\
			 &=&\dfrac{1}{1-q^4}\cdot\dfrac{1}{u_1}\cdot\dfrac{1-(q)^{20}}{\left(1-q\right)q^{19}}\\
			&=&\dfrac{1-2^{20}}{15\cdot  2^{19}}.
		\end{eqnarray*}
	} 
\end{ex}

% nguồn: Ok_Toan11_TracNghiem_Lan2
\begin{ex}%[1D2V3-6]%[ma:02eb38f8]
	Giá trị của tổng $4 + 44 + 444  + \ldots + \underbrace{44\ldots4}_{2\,018\text{ số }4}$ bằng
	\choice 
		{\True  $\dfrac{4}{9}\left(\dfrac{10^{2\,019} - 10}{9} - 2\,018\right)$}
		{$\dfrac{40}{9}(10^{2\,018} - 1) +  2\,018$}
		 {$\dfrac{4}{9}\left(10^{2\,018} - 1\right)$}
		{$\dfrac{4}{9}\left(\dfrac{10^{2\,018} -  10}{9} - 2\,018\right)$}
	 \loigiai{
		Đặt $S = 4 + 44 + 444  + \ldots + \underbrace{44\ldots4}_{2\,018\text{ số }4}$.\\
		Ta  có
		\begin{align*}
			 \dfrac{9}{4}S &= 9 + 99 + 999 + \ldots + \underbrace{99\ldots9}_{2\,018\text{ số }9}\\
			&= (10 - 1) + (10^2 - 1) + (10^3 - 1)  + \ldots + (10^{2\,018} - 1)
		\end{align*}
		Suy ra: 
		\begin{align*}
			\dfrac{9}{4}S &= (10  + 10^2 + 10^3 + \ldots + 10^{2\,018}) - 2\,018
		\end{align*}
		Với $A = 10 + 10^2 + 10^3 + \ldots + 10^{2\,018}$ là tổng $2\,018$ số hạng của  một cấp số nhân có số hạng đầu $u_1 = 10$, công bội $q = 10$ nên ta có
		\[A = u_1 \cdot \dfrac{1 - q^{2\,018}}{1 - q} = 10 \cdot \dfrac{10^{2\,018} - 1}{9} =  \dfrac{10^{2\,019} - 10}{9}.\]
		 Do đó
		\[\dfrac{9}{4}S = \dfrac{10^{2\,019} - 10}{9} - 2\,018 \Rightarrow S = \dfrac{4}{9}\left(\dfrac{10^{2\,019} -  10}{9} - 2\,018\right).\]
	} 
\end{ex}

% nguồn: Ok_Toan11_TracNghiem_Lan2
\begin{ex}%[1D2V3-6]%[ma:759d2ea3]
	Giá trị  của tổng $4+44+444+\cdots+44\ldots 4$ (với $44\ldots 4$ có $2\,018$ số $4$) bằng
	\choice 
		 {$\dfrac{40}{9} \left(10^{2\,018}-1\right)+2\,018$}
		{\True $\dfrac{4}{9} \left(\dfrac{10^{2\,019}-10}{9}-2\,018\right)$} 
		{$\dfrac{4}{9} \left(\dfrac{10^{2\,019}-10}{9}+2\,018\right)$} 
		{$\dfrac{4}{9}  \left(10^{2\,018}-1\right)$}
	 \loigiai{
		Đặt $S=4+44+444+\cdots+44\ldots  4$.\\
		Ta có 
		\begin{eqnarray*}
			\dfrac{9}{4}  S&=& 9+99+999+\cdots+99\ldots 9\\
			 &=&\left(10-1\right)+\left(10^2-1\right)+\left(10^3-1\right)+\cdots\left(10^{2\,018}-1\right).
		\end{eqnarray*}
		 Suy ra $\dfrac{9}{4} S=\left(10+10^2+10^3+\cdots+10^{2\,018} \right)-2\,018=A-2\,018$.\\
		Với $A=10+10^2+10^3+\cdots+10^{2\,018}$ là tổng $2\,018$ số hạng của một cấp số nhân có số hạng đầu $u_1=10$, công bội $q=10$  nên ta có $A=u_1 \dfrac{1-q^{2\,018}}{1-q}$ $=10\dfrac{1-10^{2\,018}}{-9}$ $=\dfrac{10^{2\,019}-10}{9}$.\\
		Do đó $\dfrac{9}{4}  S=\dfrac{10^{2\,019}-10}{9}-2\,018\Leftrightarrow S=\dfrac{4}{9} \left(\dfrac{10^{2\,019}-10}{9}-2\,018\right)$.
	} 
\end{ex}

% nguồn: Ok_Toan11_TracNghiem_Lan2
\begin{ex}%[1D2V3-6]%[ma:66ebfdfd]
	Cho dãy số xác định bởi $u_1 = 1$; $u_{n+1} = \dfrac{1}{3}\left(2u_n +  \dfrac{n - 1}{n^2 + 3n + 2}\right)$, $n \in \mathbb{N}^*$. Khi đó $u_{2\,018}$ bằng
	\choice 
		{\True $\dfrac{2^{2\,016}}{3^{2\,017}} + \dfrac{1}{2\,019}$} 
		 {$\dfrac{2^{2\,018}}{3^{2\,017}} + \dfrac{1}{2\,019}$}
		 {$\dfrac{2^{2\,017}}{3^{2\,018}} + \dfrac{1}{2\,019}$}
		{$\dfrac{2^{2\,017}}{3^{2\,018}} +  \dfrac{1}{2\,019}$}
	 \loigiai{
		 Ta có
		\begin{align*}
			u_{n+1}&= \dfrac{1}{3}\left(2u_n + \dfrac{n - 1}{n^2 + 3n +  2}\right)\\
			 &= \dfrac{1}{3}\left(2u_n + \dfrac{3}{n+2} - \dfrac{2}{n+1}\right)\\
			&= \dfrac{2}{3}u_n +  \dfrac{1}{n+2} - \dfrac{2}{3} \cdot \dfrac{1}{n+1}.
		\end{align*}
		Suy ra $u_{n+1} -  \dfrac{1}{n+2}= \dfrac{2}{3}\left(u_n - \dfrac{1}{n+1}\right) \quad (1)$\\
		Đặt  $v_n = u_n - \dfrac{1}{n+1}$, từ $(1)$ ta suy ra $v_{n+1} = \dfrac{2}{3}v_n$.\\
		Do đó $(v_n)$ là cấp số nhân với $v_1 = u_1 - \dfrac{1}{2} = \dfrac{1}{2}$, công bội $q = \dfrac{2}{3}$.\\ 
		Suy  ra
		\begin{align*}
			v_n = v_1  \cdot q^{n-1} = \dfrac{1}{2}\left(\dfrac{2}{3}\right)^{n-1}
			\Rightarrow u_n = v_n + \dfrac{1}{n+1} = \dfrac{1}{2}\left(\dfrac{2}{3}\right)^{n-1} +  \dfrac{1}{n+1}.
		\end{align*}
		Vậy $u_{2\,018} = \dfrac{1}{2}\left(\dfrac{2}{3}\right)^{2\,017} + \dfrac{1}{2\,019}  = \dfrac{2^{2\,016}}{3^{2\,017}} + \dfrac{1}{2\,019}$.
	 }
\end{ex}

% nguồn: Ok_Toan11_TracNghiem_Lan2
\begin{ex}%[1D2V3-6]%[ma:bc035412]
	Cho dãy số $\left(u_n \right)$ xác định  bởi $u_1=\dfrac{1}{3}$ và $u_{n+1}=\dfrac{n+1}{3n}\cdot u_n$. Tổng $S=u_1+\dfrac{u_2}{2}+\dfrac{u_3}{3}+\cdots+\dfrac{u_{10}}{10}$ bằng
	\choice 
		 {$\dfrac{3\,280}{6\,561}$}
		{\True $\dfrac{29\,524}{59\,049}$} 
		{$\dfrac{25\,942}{59\,049}$} 
		{$\dfrac{1}{243}$} 
	 \loigiai{
		Theo đề ta có $u_{n+1}=\dfrac{n+1}{3n}\cdot u_n$  $\Leftrightarrow \dfrac{u_{n+1}}{n+1}=\dfrac{1}{3} \dfrac{u_n}{n}$ mà $u_1=\dfrac{1}{3}$ hay $\dfrac{u_1}{1}=\dfrac{1}{3}$.\\
		Nên ta có $\dfrac{u_2}{2}=\dfrac{1}{3}\cdot \dfrac{1}{3}=\left(\dfrac{1}{3}\right)^2$; $\dfrac{u_3}{3}=\dfrac{1}{3}\cdot \left(\dfrac{1}{3} \right)^2=\left(\dfrac{1}{3}\right)^3$; $\ldots$; $\dfrac{u_{10}}{10}=\left(\dfrac{1}{3} \right)^{10}$.\\ 
		Hay dãy $\left(\dfrac{u_n}{n} \right)$ là một cấp số nhân có số hạng đầu $u_1=\dfrac{1}{3}$, công bội  $q=\dfrac{1}{3}$.\\
		 Khi đó $S=u_1+\dfrac{u_2}{2}+\dfrac{u_3}{3}+\cdots+\dfrac{u_{10}}{10}=\dfrac{3^{10}-1}{2\cdot 3^{10}} =\dfrac{59\,048}{2\cdot 3^{10}} =\dfrac{29\,524}{59\,049}$.
	 }
\end{ex}

% nguồn: Ok_Toan11_TracNghiem_Lan2
\begin{ex}%[1D2C3-6]%[ma:1d66d715]
	Số thập phân vô hạn tuần hoàn $3{,}1\,555\ldots=3{,}1(5)$ viết dưới dạng  số hữu tỉ là
	\choice 
		 {$\dfrac{63}{20}$}
		{\True  $\dfrac{142}{45}$}
		 {$\dfrac{1}{18}$}
		 {$\dfrac{7}{2}$}
	\loigiai{ 
		 $3{,}1\,555\ldots=3{,}1+0{,}05+0{,}005+0{,}0\,005+\ldots$.\\
		Dãy số  $0{,}05$; $0{,}005$; $0{,}0\,005$; $0{,}00\,005$; $\ldots$ là một cấp số nhân lùi vô hạn có $u_1=0{,}05$; $q=0{,}1$.\\
		 Vậy $3{,}1\,555\ldots=3{,}1+\dfrac{0{,}05}{1-0{,}1}=\dfrac{142}{45}$.
	} 
\end{ex}

% nguồn: Ok_Toan11_TracNghiem_Lan2
\begin{ex}%[1D2C3-6]%[ma:4bd77e69]
	Tính tổng $S=-1+\dfrac{1}{6}-\dfrac{1}{6^2}+\cdots  +\left(-1\right)^{n-1}\dfrac{1}{6^n}+\cdots$.
	 \choice
		{$S=\dfrac{7}{6}$} 
		 {\True $S=-\dfrac{6}{7}$}
		 {$S=\dfrac{6}{7}$}
		{$S=-\dfrac{7}{6}$} 
	\loigiai{ 
		Ta  có $q=\dfrac{u_2}{u_1}=\dfrac{u_3}{u_2}=\ldots=-\dfrac{1}{6}\left(|q|<1\right)$.\\
		Do  đó $S=\dfrac{u_1}{1-q}=\dfrac{-1}{1+\dfrac{1}{6}}=-\dfrac{6}{7}$.
	 }
\end{ex}


%%%%%%%%%%%%%%%% ĐÚNG/SAI %%%%%%%%%%%%%%%%

% nguồn: Ok_dungsai_lop11_Lan1
\begin{ex}%[1D2H3-6]%[ma:c221f556]
 Cho cấp số nhân ${\left(u_n\right)}$ thoả mãn $\heva{&u_4+u_6=-540 \\& u_3+u_5=180}$. Khi đó: 
  \choiceTF
  {\True Số hạng $u_1=2$} 
  {Gọi $q$ là công bội của cấp số nhân, thì ba số  $q;1;3$ tạo thành một cấp số cộng}
  {Số $-486$ là số hạng thứ $5$ của cấp số nhân } 
  {\True Tổng của $21$ số hạng đầu cấp số nhân đã cho bằng  $5230176602$}
 \loigiai{ 
  Gọi ${q}$ là công bội và ${S_{21}}$ là  tổng của 21 số hạng đầu của cấp số nhân ${\left(u_n\right)}$.
   Ta có: \allowdisplaybreaks
  \begin{eqnarray*}
   \heva{&u_4+u_6=-540  \\ &u_3+u_5=180} &\Leftrightarrow&\heva{&\left(u_3+u_5\right) q=-540 \\ &u_3+u_5=180}\\ &\Leftrightarrow&\heva{&180 q=-540 \\ &u_3+u_5=180}\\
   &\Leftrightarrow&\heva{&q=-3 \\&u_1(-3)^2+u_1(-3)^4=180}\\ 
   &\Leftrightarrow&\heva{& q = - 3 \\&u _  { 1 } ( 9 + 8 1) = 1 8 0 }\\
   &\Leftrightarrow  &\heva{&q=-3 \\&u_1=2.}\\
  \end{eqnarray*}
  Số $-486=2\cdot (-3)^5$ nên số $-486$ là số hạng  thứ $6$.\\
  Suy ra ${S_{21}=\dfrac{u_1\left(1-q^{21}\right)}{1-q}=\dfrac{2\left[1-(-3)^{21}\right]}{1-(-3)}=\dfrac{1+3^{21}}{2}}$. 
  }
\end{ex}

% nguồn: Ok_dungsai_lop11_Lan1
\begin{ex}%[1D2H3-6]%[ma:19ced7c9]
 Cho cấp số nhân ${\left(u_n\right)}$  thoả mãn $\heva{&u_4=\dfrac{2}{27} \\ &u_3=243 u_8}$. Khi đó
  \choiceTF
  {Số hạng $u_1=2;\, u_2=\dfrac{2}{3};\,  u_3=\dfrac{2}{9};\, u_4=\dfrac{2}{27};\, u_5=\dfrac{2}{81}$}
  { $u_5-u_3=-\dfrac{16}{81}$} 
  {Số ${\dfrac{2}{6561}}$ là  số hạng thứ $8$ của cấp số nhân}
  {Tổng  chín số hạng đầu của cấp số nhân là số lớn hơn $3$}
  \loigiai{
  \begin{enumerate}
   \item Gọi $q$ là công bội của cấp số  nhân $\left(u_n\right)$.\\
   Theo giả thiết, ta có $\heva{&u_1 q^3=\dfrac{2}{27} \\ &u_1 q^2=243 \cdot u_1 q^7} \Rightarrow \heva{&u_1  q^3=\dfrac{2}{27} \\ &q^5=\dfrac{1}{243}} \Rightarrow \heva{&q=\dfrac{1}{3} \\ &u_1=2}$.\\
   Năm số hạng đầu của $\left(u_n\right)$ là $u_1=2;\, u_2=\dfrac{2}{3};\, u_3=\dfrac{2}{9};\, u_4=\dfrac{2}{27};\, u_5=\dfrac{2}{81}$. 
   \setcounter{enumi}{2} 
   \item Số hạng tổng  quát của cấp số nhân $u_n=u_1 q^{n-1}=\dfrac{2}{3^{n-1}}$.\\
    Xét $u_n=\dfrac{2}{6561} \Rightarrow \dfrac{2}{3^{n-1}}=\dfrac{2}{6561}\Rightarrow 3^{n-1}=6561=3^8 \Rightarrow n=9$.\\
   Vậy $\dfrac{2}{6561}$ là số hạng thứ $9$ của cấp  số nhân $\left(u_n\right)$.
   \item Tổng chín số hạng đầu  của cấp số nhân là $$S_9=\dfrac{u_1\left( 1-{{q}^{9}} \right)}{1-q}=\dfrac{2\cdot\left( 1-\left( \dfrac{1}{3} \right)^9 \right)}{1-\dfrac{1}{3}}\approx 2{,}99985<3.$$
  \end{enumerate}
 } 
\end{ex}

% nguồn: Ok_dungsai_lop11_Lan1
\begin{ex}%[1D2H3-6]%[ma:f822020c]
 Cho cấp số nhân $\left(u_n\right)$ có tổng $n$ số hạng đầu tiên là $S_n=5^n-1$ với $n=1,2,\ldots$. 
 \choiceTF 
  {Công  bội của cấp số nhân đã cho là $q=6$ }
  { \True Số hạng đầu của  cấp số nhân đã cho là $u_1=4$}
  {\True  Tổng $5$ số hạng đầu của cấp số nhân là $ 3124$}
  {Số hạng thứ $5$ của  cấp số nhân là $ 2600$}
 \loigiai{ 
  \begin{itemchoice}
   \itemch Ta  có $\heva{&{u_1}=S_1=5-1=4\\&{u_1}+u_2=S_2=5^2-1=24
   }\Rightarrow  \heva{&{u_1}=4\\&{u_2}=24-u_1=20
    .}$\\
   Do đó $u_1=4$, và $q=\dfrac{u_2}{u_1}=5$. 
    \itemch $u_1=4$.
   \itemch  $S_5=5^5-1=3124$.
   \itemch Ta  có $u_5=u_1.q^4=4\cdot5^4=2500$.
  \end{itemchoice}
 } 
\end{ex}

% nguồn: Ok_dungsai_lop11_Lan1
\begin{ex}%[1D2H3-6]%[ma:3dc177db]
 Cho cấp  số nhân ${\left(u_n\right)}$ biết rằng ${u_1+u_2+u_3=168}$ và ${u_4+u_5+u_6=21}$. Khi đó:
  \choiceTF
  {Số hạng $u_1=90$} 
  {Công bội của cấp số  nhân bằng $2$}
  {\True Số  $24$ là số hạng thứ 3 của cấp số nhân }
  {\True Tổng của 10 số hạng đầu cấp  số nhân đã cho bằng $\dfrac{3069}{16}$}
 \loigiai{ 
  Ta  có: \allowdisplaybreaks
  \begin{eqnarray*}
   \heva{&u_1+u_2+u_3=168  \\ &u_4+u_5+u_6=21}
   &\Leftrightarrow&\heva{&u_1+u_1 \cdot q+u_1 \cdot q^2=168 \\ &u_1 \cdot q^3+u_1 \cdot  q^4+u_1 \cdot q^5=21}\\
   &\Leftrightarrow&\heva{& u_1 ( 1 + q + q^2)  = 168 \\& u _1 q^3 (1 + q + q^2) = 21} \\
   &\Leftrightarrow& \heva{&u_1 = \dfrac{ 168 }{ 1 + q + q^ 2} \\&q ^ {  3 } = \dfrac{ 1 }{ 8 }}\\
    & \Leftrightarrow& \heva{&u_1=96 \\&q=\dfrac{1}{2}}
  \end{eqnarray*}
  Ta  có $24=96\cdot \left( \dfrac{1}{2} \right)^{3-1}$.\\
  Ta  có $S_{10}=\dfrac{u_1\left( 1-q^{10} \right)}{1-q}=\dfrac{96\left[1-\left( \dfrac{1}{2}\right) ^{10} \right]}{1-\dfrac{1}{2}}=\dfrac{3069}{16}.$
 } 
\end{ex}

% nguồn: Ok_dungsai_lop11_Lan1
\begin{ex}%[1D2H3-6]%[ma:13ae0bf1]
 Cho cấp số nhân $\left(u_n\right)$ có $u_1=2$; $u_2=-4$. Xác định tính  đúng, sai của các phát biểu sau đây
 \choiceTF 
   {Công bội $ q=2$}
  {$u_5=-32$} 
  {\True Số $-64$ là số hạng thứ  $6$ của $\left(u_n\right)$}
  {\True Tổng của $8$ số hạng đầu tiên của cấp số nhân bằng $-170$} 
 \loigiai{ 
  \begin{itemchoice}
   \itemch  $ q=\dfrac{u_2}{u_1}=\dfrac{-4}{2}=-2$.
   \itemch $u_5=u_1\cdot q^4=2\cdot\left(-2\right)^4=32$. 
   \itemch $u_6=u_1\cdot q^5=2\cdot\left(-2\right)^5=-64$. 
   \itemch  $S_8=\dfrac{u_1\left(1-q^8\right)}{1-q}=\dfrac{2\left[1-\left(-2\right)^8\right]}{1-\left(-2\right)}=-170$.
  \end{itemchoice}
 } 
\end{ex}

% nguồn: Ok_dungsai_lop11_Lan2
\begin{ex}%[1D2H3-6]%[ma:d6801804]
	Cho cấp số nhân $\left(u_n\right)$ có  $u_2=6$ và $u_3=18$. Các mệnh đề sau đúng hay sai?
	 \choiceTF
		{\True Công bội của cấp số  nhân là $q=3$}
		{$u_{k-1}+u_{k+1}=u_k^2$, $2 \leq k \in  \mathbb{N}^*$}
		{Số hạng đầu của cấp  số nhân là $u_1=1$}
		{\True Tổng $10$ số hạng đầu tiên của câp  số nhân đã cho $S_{10}=59\,048$}
	 \loigiai{
		\begin{itemchoice}
			\itemch  \textbf{Đúng} \\
			Công bội của  cấp số nhân là $q=\dfrac{u_3}{u_2}=3$.
			\itemch  \textbf{Sai} \\
			Do $\left(u_n\right)$ là cấp số nhân $u_{k-1} \cdot u_{k+1}=u_k^2$, $2 \leq k \in  \mathbb{N}^*$.
			 \itemch \textbf{Sai} \\
			Số hạng đầu của  cấp số nhân là $u_1=\dfrac{u_2}{q}=2$.
			 \itemch \textbf{Đúng} \\
			Tổng $10$ số hạng  đầu tiên của câp số nhân đã cho $S_{10}=2\cdot\dfrac{1-3^{10}}{1-3}=59\,048$.
		\end{itemchoice}
	} 
\end{ex}

% nguồn: Ok_dungsai_lop11_Lan2
\begin{ex}%[1D2H3-6]%[ma:05475863]
	Cho cấp số nhân $\left(u_n \right)$ thoả  mãn: $\heva{&{u_4=\dfrac{2}{27}} \\
	&{u_3=243u_8} 
	}$. Khi  đó
	 \choiceTF
		{\True Số  hạng $u_1=2;u_2=\dfrac{2}{3}$}
		{\True  $u_5-u_3=-\dfrac{16}{81}$}
		{Số $\dfrac{2}{6\,561}$ là số  hạng thứ $8$ của cấp số nhân}
		{Tổng chín  số hạng đầu của cấp số nhân là số lớn hơn $3$}
	 \loigiai{
		\begin{itemchoice}
			\itemch \textbf{Đúng.} \\ 
			Gọi $q$ là công bội  của cấp số nhân $\left(u_n \right)$.\\
			Theo  giả thiết, ta có
			\[\heva{&u_1q^3=\dfrac{2}{27}\\& u_1q^2=243\cdot u_1q^7} \Rightarrow \heva{&u_1q^3=\dfrac{2}{27}\\ &q^5=\dfrac{1}{243}} \Rightarrow  \heva{&q=\dfrac{1}{3}\\&u_1=2.}\]
			Năm số  hạng đầu của $\left(u_n \right)$ là $u_1=2;u_2=\dfrac{2}{3};u_3=\dfrac{2}{9};u_4=\dfrac{2}{27};u_5=\dfrac{2}{81}$
			\itemch  \textbf{Đúng.} \\ $u_5-u_3=-\dfrac{16}{81}$.
			 \itemch \textbf{Sai.} \\
			Số  hạng tổng quát của cấp số nhân: $u_n=u_1 q^{n-1}=\dfrac{2}{3^{n-1}}$.
			 Xét $u_n=\dfrac{2}{6561} \Rightarrow \dfrac{2}{3^{n-1}}=\dfrac{2}{6561}$
			 \[\Rightarrow 3^{n-1}=6561=3^8\Rightarrow n=9.
			 \]
			Vậy $\dfrac{2}{6561}$ là số hạng thứ 9 của cấp  số nhân $\left(u_n \right)$.
			\itemch \textbf{Sai.} \\ 
			Tổng chín số hạng đầu của cấp số nhân là $S_9=\dfrac{u_1 \left(1-q^9\right)}{1-q}=\dfrac{2.\left(1-\left(\dfrac{1}{3} \right)^9\right)}{1-\dfrac{1}{3}} \approx 2{,}99985< 3$. 
		\end{itemchoice}
	 }
\end{ex}

% nguồn: Ok_dungsai_lop11_Lan2
\begin{ex}%[1D2H3-6]%[ma:94d3d492]
	Cho dãy số $\left(u_n\right)$ là một cấp số nhân có $20$ số hạng, biết số hạng đầu của cấp số nhân đó là  $2$ và số hạng thứ $201$ là $-2^{20}$. Các mệnh đề sau đây đúng hay sai?
	\choiceTF 
		{Công bội của cấp số nhân đó  là $q=2$}
		{Số hạng thứ $10$ của cấp số nhân đó  bằng đó $2^{10}$}
		{Tổng $18$  số hạng đầu của cấp số nhân đó bằng $\dfrac{2}{3} \cdot\left(1+2^{18}\right)$}
		 {\True $\left(u_5\right)^2=u_4 \cdot u_6$}
	 \loigiai{
		\begin{itemchoice}
			\itemch \textbf{Sai} \\ 
			Vì dãy số $\left(u_n\right)$ là một cấp  số nhân có $u_1=2$, $u_{20}=-2^{20}$ suy ra
			\[u_{20}=u_1 \cdot q^{19} \Rightarrow 2 \cdot q^{19}=-2^{20}  \Rightarrow q^{19}=(-2)^{19} \Rightarrow q=-2.\]
			\itemch \textbf{Sai} \\ 
			Số  hạng thứ $10$ của cấp số nhân đó là $u_{10}=u_1 \cdot q^9=2 \cdot(-2)^9=-2^{10}$.
			 \itemch \textbf{Sai} \\
			Tổng của $19$ số hạng đầu của cấp số nhân đó bằng $S_{19}=\dfrac{2 \cdot\left[1-(-2)^{18}\right]}{1-(-2)}=\dfrac{2}{3} \cdot\left(1-2^{18}\right)$. 
			\itemch \textbf{Đúng} \\ 
			 Ta có $u_4 \cdot u_6=\dfrac{u_5}{q} \cdot u_5 \cdot q=u_5 \cdot u_5=\left(u_5\right)^2$.
		\end{itemchoice}
	} 
\end{ex}

% nguồn: Ok_dungsai_lop11_Lan2
\begin{ex}%[1D2H3-6]%[ma:ebdc3948]
	Cho cấp số  nhân $\left(u_n \right)$ có công bội nguyên và các số hạng thoả mãn $\heva{& u_4-u_2=54 \\ &u_5-u_3=108.}$
	\choiceTF 
		{\True Số hạng đầu của cấp số nhân bằng $9$} 
		{Công bội của cấp số nhân $q=3$} 
		{\True Tổng của $9$ số hạng đầu tiên  bằng $4\,599$}
		{Số $576$ là số hạng thứ 6  của cấp số nhân}
	 \loigiai{
		\begin{itemchoice}
			\itemch  \textbf{Đúng.} \\
			Ta  có
			\begin{eqnarray*}
				 \heva{&u_4-u_2=54\\
				&u_5-u_3=108}  &\Leftrightarrow &\heva{&u_1q^3-u_1q=54\\
				&u_1q^4-u_1q^2=108}  \\
				 &\Leftrightarrow& \heva{&u_1q\left(q^2-1\right)=54\\
				&u_1q^2\left(q^2-1\right)=108}\\ 
				&\Leftrightarrow& \heva{&u_1=\dfrac{54}{q  (q ^2-1)} \\&
				\dfrac{1}{q}=\dfrac{54}{108}}\\ 
				&\Leftrightarrow& \heva{&u_1=\dfrac {54}{2(2^2-1)} \\&q=2}\\ &\Leftrightarrow& \heva{&u_1=9\\&q=2.} 
			\end{eqnarray*}
			\itemch \textbf{Sai.}  \\
			Công bội của cấp số nhân  $q=9$.
			\itemch  \textbf{Đúng.} \\
			Ta có $S_n=4\,599\Leftrightarrow \dfrac{u_1 \cdot \left(1-q^n \right)}{1-q}=4\,599\Leftrightarrow  \dfrac{9\cdot \left(1-2^n \right)}{1-2}=4\,599$
			\[\Leftrightarrow-9\cdot  \left(1-2^n \right)=4\,599\Leftrightarrow 1-2^n=-511\Leftrightarrow 2^n=512\Leftrightarrow n=9.\]
			Vậy tổng của $9$ số hạng đầu  tiên bằng $4\,599$.
			\itemch  \textbf{Sai.} \\
			 Ta có
			\begin{eqnarray*}
				 u_k=576 &\Leftrightarrow& u_1 \cdot q^{k-1}=576\\
				&\Leftrightarrow& 9\cdot 2^{k-1}=576\\ 
				 &\Leftrightarrow& 2^{k-1}=64\\
				&\Leftrightarrow& k-1=6\Leftrightarrow  k=7.
			\end{eqnarray*}
			Vậy số $576$  là số hạng thứ $7$ của cấp số nhân.
		\end{itemchoice}
	 }
\end{ex}

% nguồn: Ok_dungsai_lop11_Lan2
\begin{ex}%[1D2H3-6]%[ma:07189f67]
	Cho dãy số $\left(u_n\right)$, biết $\heva{&u_1=3 \\ &u_{n+1}=4u_n-1}$, $\forall n \in \mathbb{N}^*$. Các mệnh đề sau đúng hay sai? 
	\choiceTF 
		{Số hạng thứ năm  của đãy số là $685$}
		{\True Đặt $v_n=u_n-\dfrac{1}{3}$. Ta  có $\left(v_n\right)$ là cấp số nhân}
		{\True Số hạng tổng  quát $u_n=\dfrac{8}{3} \cdot 4^{n-1}+\dfrac{1}{3}$}
		 {\True Ta có $S_8=58\,256$}
	\loigiai{ 
		\begin{itemchoice}
			\itemch \textbf{Sai} \\ 
			Ta có 
			\begin{eqnarray*}
				&&u_2=4u_1-1=4\cdot  3-1=12-1=11\\
				 &&u_3=4u_2-1=4\cdot 11-1=44-1=43\\
				 &&u_4=4u_3-1=4\cdot 43-1=172-1=171\\
				&&u_5=4u_4-1=4\cdot 171-1=684-1=683. 
			\end{eqnarray*}
			 Vậy số hạng thứ năm của dãy số là $u_5=683$.
			 \itemch \textbf{Đúng} \\
			Ta có $u_{n+1}=4u_n-1  \Leftrightarrow u_{n+1}-\dfrac{1}{3}=4\left(u_n-\dfrac{1}{3}\right)$.\\
			Đăt $v_n=u_n-\dfrac{1}{3}$. Ta có  $v_{n+1}=4v_n$. Suy ra $\left(v_n\right)$ là cấp số nhân.
			 \itemch \textbf{Đúng} \\
			Ta có $v_{n+1}=4 v_n$. Suy ra $\left(v_n\right)$ là cấp  số nhân với $\heva{&q=4 \\& v_1=u_1-\dfrac{1}{3}=\dfrac{8}{3}.}$\\
			Suy ra $v_n=v_1 \cdot q^{n-1}=\dfrac{8}{3} \cdot 4^{n-1} \Rightarrow u_n=\dfrac{8}{3} \cdot 4^{n-1}+\dfrac{1}{3}$. 
			\itemch  \textbf{Đúng} \\
			Ta  có $S_8=u_1+u_2+\ldots+u_8=\dfrac{8}{3} \cdot\left(1+4+4^2+\ldots+4^7\right)+\dfrac{1}{3} \cdot 8=58\,256$.
		\end{itemchoice}
	 }
\end{ex}

% nguồn: Ok_dungsai_lop11_Lan2
\begin{ex}%[1D2H3-6]%[ma:536e65e4]
	Xác định  tính đúng, sai của các khẳng định sau
	 \choiceTF
		{\True Dãy số $\left(u_n \right)$ với $u_n=(-7)^n \cdot 5^{3n-1}$  là cấp số nhân với công bội $q=-875$}
		{Dãy số $\left(u_n \right)$ với $\heva{&u_1=2 \\ &u_{n+1}=-5u_n}$ là cấp  số nhân với công bội $q=-4$}
		{\True Dãy số $\left(u_n \right)$ với $\heva{&u_1=2 \\ &u_{n+1}=u_n^2} $ không là cấp số  nhân}
		{\True Dãy số $\left(u_n  \right)$ với $-\dfrac{1}{8};-\dfrac{1}{4};-\dfrac{1}{2};1$ không là cấp số nhân}
	\loigiai{ 
		\begin{itemchoice}
			\itemch  \textbf{Đúng.} \\
			Ta có $\dfrac{u_{n+1}}{u_n}=\dfrac{(-7)^{n+1} \cdot 5^{3(n+1)-1}}{(-7)^n \cdot  5^{3n-1}}=\dfrac{(-7)\cdot 5^2}{5^{-1}}=-7\cdot 5^3=-875$ không đổi.\\
			Vậy  $\left(u_n \right)$ là cấp số nhân với công bội $q=-875$.
			\itemch \textbf{Sai.} \\ 
			Ta có $\heva{&{u_1=2}  \\
			&{u_{n+1}=-5u_n} 
			}$.  \\
			Khi đó $\dfrac{u_{n+1}}{u_n}=\dfrac{-5u_n}{u_n}=-5$.\\ 
			Vậy $\left(u_n \right)$ là cấp số nhân với công  bội $q=-5$.
			 \itemch \textbf{Đúng.} \\ Ta có $u_2=u_1^2=2^2=4;u_3=u_2^2=4^2=16;u_4=u_3^2=16^2=256$.\\
			Khi  đó $\dfrac{u_2}{u_1}=\dfrac{4}{2}=2;\dfrac{u_4}{u_3}=\dfrac{256}{16}=16$.\\
			Nhận thấy $\dfrac{u_2}{u_1} \ne \dfrac{u_4}{u_3} (2\ne 16)$.\\ 
			Vậy $\left(u_n \right)$ không  là cấp số nhân.
			 \itemch \textbf{Đúng.} \\
			Ta có $\dfrac{u_2}{u_1}=\left(-\dfrac{1}{4} \right):\left(-\dfrac{1}{8} \right)=2;\dfrac{u_4}{u_3}=1:\left(-\dfrac{1}{2}  \right)=-2$.\\
			Nhận thấy: $\dfrac{u_2}{u_1} \ne \dfrac{u_4}{u_3}  (2\ne-2)$.\\
			Vậy $\left(u_n \right)$ không  là cấp số nhân.
		\end{itemchoice}
	} 
\end{ex}

% nguồn: Ok_dungsai_lop11_Lan2
\begin{ex}%[1D2H3-6]%[ma:a5c7efe3]
	Cho cấp số nhân $\left(u_n \right)$, biết $u_1+u_5=51$; $u_2+u_6=102$. Khi  đó
	 \choiceTF
		{\True  Số hạng $u_1=3$}
		{Số hạng  $u_4=48$}
		{Số $12\,288$ là  số hạng thứ $12$ của cấp số nhân $\left(u_n \right)$}
		{\True Tổng tám số hạng đầu của cấp số nhân  là $765$}
	 \loigiai{
		\begin{itemchoice}
			 \itemch \textbf{Đúng.} \\
			Gọi $q$  là công bội của cấp số nhân đã cho. Ta có
			 \[\heva{&{u_1+u_5=51} \\
			&{u_2+u_6=102}}\Leftrightarrow \heva{&{u_1+u_1q^4=51}  \\
			&{u_1q+u_1q^5=102}}\Leftrightarrow  \heva{&{u_1\left(1+q^4\right)=51 \quad\qquad (1)} \\
			&{u_1q\left(1+q^4\right)=102.\qquad  (2)} } \]
			Nhận xét: Nếu $u_1=0$ hay $q=0$  thì (1) và (2) đều không thoả mãn, vì vậy ta có $u_1 q\ne 0$. Chia theo vế (2) cho (1), ta được $q=2$.
			Thay $q=2$ vào (1) suy  ra $u_1=\dfrac{51}{1+2^4}=3$.
			Công thức số hạng tổng quát của cấp số nhân $u_n=3\cdot  2^{n-1}$
			\itemch \textbf{Sai.}\\ 
			 $u_4=3\cdot2^3=24$.
			\itemch \textbf{Sai.}\\ 
			Xét $u_n=12\,288\Leftrightarrow 3\cdot2^{n-1}=12\,288\Leftrightarrow 2^{n-1}=2^{12} \Leftrightarrow  n=13$.\\
			Vậy $12\,288$ là  số hạng thứ $13$ của cấp số nhân đã cho.
			\itemch \textbf{Đúng.}  \\
			Tổng  tám số hạng đầu của cấp số nhân là $S_8=\dfrac{u_1 \left(1-q^8\right)}{1-q}=\dfrac{3\cdot \left(1-2^8\right)}{1-2}=765$.
		\end{itemchoice}
	 }
\end{ex}

% nguồn: Ok_dungsai_lop11_Lan2
\begin{ex}%[1D2H3-6]%[ma:515a12bb]
	Cho dãy số $\left(u_n\right)$ là một cấp số nhân cón $u_1=3$, $u_3=12$ và công bội của cấp số nhân đó là số âm. Các mệnh đề sau  đúng hay sai?
	 \choiceTF
		{\True Công bội của cấp số cộng đó  là $q=-2$}
		{Số hạng thứ $25$ của cấp số nhân đó bằng $-3\cdot 2^{24}$} 
		{Tổng $101$ số hạng đầu của cấp số nhân đó bằng  $1-2^{101}$}
		{\True  $u_{55}=\sqrt{u_{54} \cdot u_{56}}$}
	\loigiai{ 
		\begin{itemchoice}
			\itemch  \textbf{Đúng} \\
			Vì dãy số  $\left(u_n\right)$ là một cấp số nhân có $u_1=3$, $u_3=12$ và công bội $q<0$, suy ra
			\[u_3=u_1  \cdot q^2 \Rightarrow 3 \cdot q^2=12 \Rightarrow q=-2.\]
			\itemch  \textbf{Sai} \\
			Ta  có $u_{25}=u_1 \cdot q^{24}=3 \cdot(-2)^{24}=3 \cdot 2^{24}$.
			\itemch \textbf{Sai}  \\
			Tổng $101$ số hạng đầu  của cấp số nhân đó là $S_{101}=\dfrac{u_1\left(1-q^{101}\right)}{1-q}=\dfrac{3 \cdot\left[1-(-2)^{101}\right]}{1-(-2)}=1+2^{101}$.
			\itemch \textbf{Đúng} \\ 
			Ta có $\sqrt{u_{54} \cdot u_{56}}=\sqrt{\dfrac{u_{55}}{q} \cdot u_{55} \cdot q}=\sqrt{\left(u_{55}\right)^2}=\left|u_{55}\right|=u_{55}$ $\left(\text{vì }u_{55}=u_1  \cdot q^{54}>0\right)$.
		\end{itemchoice}
	 }
\end{ex}

% nguồn: Ok_dungsai_lop11_Lan2
\begin{ex}%[1D2H3-6]%[ma:22472cc6]
	Cho cấp số  nhân $\left(u_n\right)$ có $u_1=-3$ và $q=-2$. Các mệnh đề sau đúng hay sai?
	 \choiceTF
		 {\True Số hạng thứ $2$ của cấp số nhân đã cho là $u_2=6$}
		{Số  $192$ là số hạng hạng thứ $7$}
		{\True Tổng $10$ số hạng đầu tiên của câp số nhân  đã cho $S_{10}=1\,023$}
		{\True Số hạng  tổng quát của cấp số nhân đã cho là $u_n=-3 \cdot(-2)^{n-1}$}
	\loigiai{ 
		\begin{itemchoice}
			\itemch  \textbf{Đúng} \\
			Số  hạng thứ $2$ của cấp số nhân đã cho là $u_2=u_2 \cdot q=(-3) \cdot(-2)=6$.
			\itemch \textbf{Sai} \\ 
			Số hạng  thứ $7$ của cấp số nhân đã cho là $u_7=u_1 \cdot q^6=(-3) \cdot(-2)^6=-192$.
			\itemch  \textbf{Đúng} \\
			Tổng  10 số hạng đầu tiên của câp số nhân đã cho $S_{10}=-3 \dfrac{1-(-2)^{10}}{1-(-2)}=1\,023$.
			 \itemch \textbf{Đúng} \\
			Số hạng tổng quát của cấp số nhân đã cho  là $u_n=-3 \cdot(-2)^{n-1}$.
		\end{itemchoice}
	 }
\end{ex}

% nguồn: Ok_dungsai_lop11_Lan2
\begin{ex}%[1D2H3-6]%[ma:07a04f3c]
	Cho cấp số nhân $\left(u_n \right)$ thoả  mãn $\heva{&u_4+u_6=-540 \\ & u_3+u_5=180} $. Khi đó
	\choiceTF 
		{\True Số  hạng $u_1=2$}
		{Gọi $q$ là công bội của cấp số nhân, thì  ba số $q;1;3$ tạo thành một cấp số cộng}
		{Số $-486$ là số  hạng thứ $5$ của cấp số nhân}
		{\True Tổng của $21$ số hạng đầu cấp số nhân đã cho bằng $5\,230\,176\,602$} 
	\loigiai{ 
		Gọi $q$ là công bội và $S_{21}$ là  tổng của $21$ số hạng đầu của cấp số nhân $\left(u_n \right)$.
		 Ta có
		\begin{eqnarray*}
			\heva{& u_4+u_6=-540 \\&u_3+u_5=180 } &\Leftrightarrow&  \heva{& \left(u_3+u_5\right)q=-540\\& u_3+u_5=180 }\\
			&\Leftrightarrow& \heva{& 180q=-540\\& u_3+u_5=180 }  \Leftrightarrow \heva{&q=-3\\& u_1(-3)^2+u_1(-3)^4=180}\\
			&\Leftrightarrow& \heva{&q=-3\\&u_1=2.} 
		\end{eqnarray*}
		\begin{itemchoice}
			\itemch \textbf{Đúng.} \\ 
			 Số hạng $u_1=2$.
			\itemch  \textbf{Sai.} \\
			Gọi $q$ là công bội của cấp số nhân, thì ba số $-3;1;3$ không tạo  thành một cấp số cộng.
			\itemch \textbf{Sai.} \\ 
			Số $-486=2\cdot  (-3)^5$ nên số $-486$ là số hạng thứ $6$.
			\itemch \textbf{Đúng.} \\ 
			Ta  có $S_{21}=\dfrac{u_1 \left(1-q^{21} \right)}{1-q}=\dfrac{2\left[1-(-3)^{21} \right]}{1-(-3)}=\dfrac{1+3^{21}}{2}$.
		\end{itemchoice}
	 }
\end{ex}

% nguồn: Ok_dungsai_lop11_Lan2
\begin{ex}%[1D2H3-6]%[ma:47da2c27]
	Một cấp số nhân có số hạng đầu $u_1=3$, công bội $q=2$. Xét tính đúng sai của các  mệnh đề sau
	\choiceTF 
		{\True Số  hạng $u_2=6$}
		{\True Số hạng tổng quát $u_n=3\cdot 2^{n-1}$ $(n\geq 1)$} 
		{\True Biết $S_n=765$ thì  $n=8$}
		 {\True Số hạng $u_{2n}>0$}
	 \loigiai{
		\begin{itemchoice}
			 \itemch \textbf{Đúng} \\
			Số hạng $u_2=u_1\cdot d=3\cdot  2=6$.
			 \itemch \textbf{Đúng} \\
			Số hạng tổng quát $u_n=u_1\cdot q^{n-1}=3\cdot 2^{n-1}$  $(n\geq 1)$.
			\itemch \textbf{Đúng}  \\
			Ta có $S_n=765\Leftrightarrow u_1\dfrac{1-q^n}{1-q}=765\Leftrightarrow 3\cdot \dfrac{1-2^n}{1-2}=765\Leftrightarrow 2^n=256\Leftrightarrow n=8$. 
			 \itemch \textbf{Đúng} \\
			Ta có $u_{2n}=3\cdot 2^{2n-1}=\dfrac{3}{2}\cdot  4^n>0$.
		\end{itemchoice}
	} 
\end{ex}

% nguồn: Ok_dungsai_lop11_Lan2
\begin{ex}%[1D2H3-6]%[ma:904c4b8e]
	Cho dãy số $\left(u_n\right)$ xác định bởi  $u_1=2$ và $u_{n+1}=\dfrac{n+1}{2n} \cdot u_n$. Xét tính đúng sai của các mệnh đề sau
	\choiceTF 
		 {\True Số hạng $u_2=2$}
		{Tỉ  số $\dfrac{u_{2\,023}}{u_{2024}}=2$}
		{Số hạng  $u_{2024}=\dfrac{1}{2^{2024}}$}
		{\True Tổng  $S=u_1+\dfrac{u_2}{2}+\dfrac{u_3}{3}+\ldots+\dfrac{u_{10}}{10}$ bằng $\dfrac{1\,023}{256}$}
	 \loigiai{
		Theo đề ta có $u_{n+1}=\dfrac{n+1}{2 n} \cdot u_n \Leftrightarrow  \dfrac{u_{n+1}}{n+1}=\dfrac{u_n}{2n}$.\\
		Đặt  $v_n=\dfrac{u_n}{n}$ thì $v_{n+1}=\dfrac{1}{2} v_n$.\\
		Do đó $\left(v_n\right)$ là một cấp số nhân với  $v_1=\dfrac{u_1}{1}=2$ và công bội $q=\dfrac{1}{2}$.\\
		Số  hạng tổng quát của dãy $\left(v_n\right)$ là $v_n=v_1 \cdot q^{n-1}=2 \cdot\left(\dfrac{1}{2}\right)^{n-1}=\left(\dfrac{1}{2}\right)^{n-2} \Rightarrow u_n=n \cdot\left(\dfrac{1}{2}\right)^{n-2}$.
		\begin{itemchoice}
			\itemch \textbf{Đúng} \\ 
			Số hạng $u_2=2\cdot  \left(\dfrac{1}{2}\right)^{2-2}=2$.
			\itemch \textbf{Sai} \\ 
			 Ta có $\dfrac{u_n}{u_{n+1}}=\dfrac{2n}{n+1} \Rightarrow \dfrac{u_{2\,023}}{u_{2024}}=\dfrac{2\cdot 2\,023}{2\,024}=\dfrac{2\,023}{1\,012}$.
			\itemch \textbf{Sai} \\ 
			Ta có $u_{2\,024}=2\,024  \cdot\left(\dfrac{1}{2}\right)^{2022}=253 \cdot\left(\dfrac{1}{2}\right)^{2019}$.
			\itemch \textbf{Đúng} \\ 
			Ta có $v_n=\dfrac{u_n}{n}=\left(\dfrac{1}{2}\right)^{n-2}$ là một cấp  số nhân với $v_1=\dfrac{u_1}{1}=2$ và công bội $q=\dfrac{1}{2}$.\\
			Khi  đó $S=u_1+\dfrac{u_2}{2}+\dfrac{u_3}{3}+\ldots+\dfrac{u_{10}}{10}=v_1\dfrac{1-q^{10}}{1-q}=2\cdot \dfrac{1-\left(\dfrac{1}{2}\right)^{10} }{1-\dfrac{1}{2}}=\dfrac{1\,023}{256}$.
		\end{itemchoice}
	 }
\end{ex}

% nguồn: Ok_dungsai_lop11_Lan2
\begin{ex}%[1D2H3-6]%[ma:7edd9488]
	Cho cấp số nhân $\left(u_n \right)$ biết rằng $u_1+u_2+u_3=168$ và $u_4+u_5+u_6=21$.  Khi đó
	\choiceTF 
		 {Số hạng $u_1=90$}
		{Công bội của cấp  số nhân bằng $2$}
		{\True Số $24$ là  số hạng thứ 3 của cấp số nhân}
		{\True Tổng  của $10$ số hạng đầu cấp số nhân đã cho bằng $\dfrac{3\,069}{16}$}
	 \loigiai{
		 Ta có
		\begin{eqnarray*}
			 \heva{&u_1+u_2+u_3=168\\&u_4+u_5+u_6=21 } & \Leftrightarrow& \heva{&u_1+u_1\cdot q+u_1\cdot q^2=168 \\& u_1\cdot q^3+u_1\cdot q^4+u_1\cdot q^5=21}\\
			 & \Leftrightarrow& \heva{& u_1\left(1+q+q^2\right)=168\\&u_1q^3\left(1+q+q^2\right)=21 }\\
			& \Leftrightarrow& \heva{& u_1=\dfrac{168}{1+q+q^2}\\&  q=\dfrac{1}{2} } \Leftrightarrow \heva{&u_1=96\\&q=\dfrac{1}{2}. }
		\end{eqnarray*}
		\begin{itemchoice}
			 \itemch \textbf{Sai.} \\
			 Số hạng $u_1=96$.
			\itemch \textbf{Sai.} \\ 
			Công bội của cấp số nhân bằng $\dfrac{1}{2}$. 
			 \itemch \textbf{Đúng.}\\
			Ta có $24=96\cdot  \left(\dfrac{1}{2} \right)^{3-1}$.
			 \itemch \textbf{Đúng.}\\
			Ta có  $S_{10}=\dfrac{u_1 \left(1-q^{10} \right)}{1-q}=\dfrac{96\left[1-\left(\dfrac{1}{2}\right)^{10} \right]}{1-\dfrac{1}{2}}=\dfrac{3069}{16}$.
		\end{itemchoice}
	} 
\end{ex}

% nguồn: Ok_dungsai_lop11_Lan2
\begin{ex}%[1D2V3-6]%[ma:9663000e]
	Cho cấp số  nhân $(u_n)$, biết $u_1+u_5=51$; $u_2+u_6=102$.
	\choiceTF 
		{\True Số hạng $u_1=3$} 
		{Số  hạng $u_4=48$}
		{Số $12\,288$ là  số hạng thứ $12$ của cấp số nhân $(u_n)$}
		{\True Tổng  $8$ số hạng đầu của cấp số nhân là $765$}
	\loigiai{ 
		\begin{itemchoice}
			\itemch Gọi $q$ là công  bội của cấp số nhân đã cho.\\
			Ta có $\heva{& u_1+u_5=51  \\& u_2+u_6=102} \Leftrightarrow \heva{& u_1+u_1q^4=51 \\& u_1q+u_1q^5=102} \Leftrightarrow \heva{& u_1\left(1+q^4\right)=51 \quad (1) \\& u_1q\left(1+q^4\right)=102 \quad (2).}$\\
			Nếu $u_1=0$ hay $q=0$  thì $(1)$ và $(2)$ đều không thoả mãn vì vậy ta có $u_1q \neq 0$.\\
			Chia theo vế $(2)$ cho $(1)$ ta  được $q=2$.\\
			Thay $q=2$ vào $(1)$  suy ra $u_1 = \dfrac{51}{1+2^4}=3$.\\
			Công thức số hạng tổng quát của cấp số nhân là $u_n=3\cdot2^{n-1}$. 
			\itemch  $u_4 = 3\cdot2^3 = 24$.
			\itemch  Xét $u_n=12\,288 \Leftrightarrow 3\cdot2^{n-1}=12\,288 \Leftrightarrow 2^{n-1}=2^{12} \Leftrightarrow n=13$.\\
			Vậy $12\,288$ là số hạng thứ  $13$ của cấp số nhân đã cho.
			\itemch Tổng $8$ số hạng đầu  của cấp số nhân là $S_8 = \dfrac{u_1\left(1-q^8\right)}{1-q} = \dfrac{3\cdot\left(1-2^8\right)}{1-2}=765$.
		\end{itemchoice}
	} 
\end{ex}

% nguồn: Ok_dungsai_lop11_Lan2
\begin{ex}%[1D2V3-6]%[ma:23e0ab44]
	Cho cấp số nhân $\left(u_n\right)$ thoả mãn $\heva{& u_4=\dfrac{2}{27}  \\& u_3=243 u_8.}$
	 \choiceTF
		{\True Số hạng $u_1=2$;  $u_2=\dfrac{2}{3}$}
		{\True  $u_5-u_3=-\dfrac{16}{81}$}
		{Số $\dfrac{2}{6\,561}$  là số hạng thứ $8$ của cấp số nhân}
		{Tổng $9$ số hạng đầu  của cấp số nhân là số lớn hơn $3$}
	\loigiai{ 
		\begin{itemchoice}
			\itemch  Gọi $q$ là công bội của cấp số nhân $\left(u_n\right)$.\\
			Ta có $\heva{& u_1 q^3 = \dfrac{2}{27} \\& u_1 q^2  = 243\cdot u_1 q^7} \Leftrightarrow \heva{& u_1 q^3 = \dfrac{2}{27} \\& q^5 = \dfrac{1}{243}} \Leftrightarrow \heva{& q = \dfrac{1}{3} \\& u_1 = 2.}$\\
			Năm số hạng đầu của $\left(u_n\right)$ là $u_1  = 2$; $u_2 = \dfrac{2}{3}$; $u_3 = \dfrac{2}{9}$; $u_4 = \dfrac{2}{27}$; $u_5 = \dfrac{2}{81}$.
			 \itemch Ta có $u_5 - u_3 = \dfrac{2}{81} - \dfrac{2}{9} = -\dfrac{16}{81}$.
			\itemch Số hạng tổng quát của cấp  số nhân $u_n = u_1q^{n-1} = \dfrac{2}{3^{n-1}}.$\\
			Xét $u_n = \dfrac{2}{6\,561} \Rightarrow \dfrac{2}{3^{n-1}} = \dfrac{2}{6\,561} \Rightarrow 3^{n-1} = 6\,561 = 3^8 \Rightarrow n=9$.\\ 
			Vậy $\dfrac{2}{6\,561}$ là số hạng thứ $9$  của cấp số nhân $\left(u_n\right).$
			\itemch  Tổng $9$ số hạng đầu của cấp số nhân là $S_9 = \dfrac{u_1\left(1-q^9\right)}{1-q} = \dfrac{2\left(1-\left(\tfrac{1}{3}\right)^9\right)}{1-\tfrac{1}{3}} \approx 2{,}99985 < 3.$
		\end{itemchoice}
	 }
\end{ex}

% nguồn: Ok_dungsai_lop11_Lan2
\begin{ex}%[1D2V3-6]%[ma:a410d77e]
	Cho cấp số nhân $\left(u_n \right)$ có công bội nguyên và các  số hạng thoả mãn $\heva{& u_4-u_2=54 \\& u_5-u_3=108.}$
	 \choiceTF
		{\True Số hạng đầu của  cấp số nhân bằng $9$}
		{Công bội của cấp  số nhân $q=3$}
		 {\True Tổng của $9$ số hạng đầu tiên bằng $4\,599$}
		{Số $576$ là số hạng thứ $6$ của cấp  số nhân}
	\loigiai{Ta có $\heva{& u_4 - u_2 = 54 \\& u_5 - u_3 = 108}  \Leftrightarrow \heva{& u_1 q^3 - u_1 q = 54 \\& u_1 q^4 - u_1 q^2 = 108} \Leftrightarrow \heva{&u_1 q \left(q^2 - 1 \right) = 54 \\& u_1 q^2 \left(q^2 - 1 \right) = 108.}$ \\
		Suy ra $\heva{& u_1 = \dfrac{54}{q \left(q^2 - 1 \right) } \\& \dfrac{1}{q} = \dfrac{54}{108}} \Leftrightarrow \heva{&u_1 = \dfrac{54}{2  \left(2^2 - 1 \right) } \\& q = 2} \Leftrightarrow \heva{& u_1 = 9 \\& q = 2.}$
		\begin{itemchoice}
			\itemch Số hạng đầu của cấp số  nhân bằng $9$.
			\itemch Công  bội của cấp số nhân $q=3$.
			\itemch Ta có 
			$$\begin{aligned}
			S_n = 4\,599 
			& \Leftrightarrow u_1 \cdot \dfrac{\left(1 -  q^n \right) }{1 - q} = 4\,599 \\
			&  \Leftrightarrow \dfrac{9\cdot\left(1 - 2^n \right) }{1 - 2} = 4\,599 \\
			& \Leftrightarrow -9\cdot\left(1 - 2^n \right) = 4\,599  \\
			& \Leftrightarrow  1 - 2^n = -511 \\
			& \Leftrightarrow  2^n = 512 \\
			& \Leftrightarrow n =  9.
		\end{aligned}$$
		Vậy  tổng của $9$ số hạng đầu tiên bằng $4\,599$.
		\itemch Ta có $u_k = 576 \Leftrightarrow u_1\cdot q^{k-1} = 576 \Leftrightarrow 9\cdot 2^{k-1} = 576 \Leftrightarrow 2^{k-1} = 64 \Leftrightarrow k - 1 = 6 \Leftrightarrow k  = 7$.\\
		Vậy số $576$ là số hạng thứ $7$ của cấp số nhân. 
	\end{itemchoice}
 }
\end{ex}

% nguồn: Ok_dungsai_lop11_Lan2
\begin{ex}%[1D2V3-6]%[ma:75feac92]
	Cho cấp số nhân $(u_n)$ biết rằng $u_1+u_2+u_3=168$ và $u_4+u_5+u_6=21$. 
	 \choiceTF
		 {Số hạng $u_1=90$}
		{Công bội của cấp số nhân  bằng $2$}
		{\True Số $24$ là số  hạng thứ $3$ của cấp số nhân}
		{\True  Tổng của $10$ số hạng đầu cấp số nhân đã cho bằng $\dfrac{3\,069}{16}$}
	 \loigiai{
		\begin{itemchoice}
			\itemch Số hạng $u_1 =  96$.
			\itemch Công bội của cấp  số nhân bằng $q = \dfrac{1}{2}$.
			\itemch Ta có  $24 = 96\cdot\left(\dfrac{1}{2} \right)^{3-1}$.
			\itemch Ta có $S_{10} = \dfrac{u_1(1-q^{10})}{1-q} =  \dfrac{96\left[1-\left(\tfrac{1}{2}\right)^{10} \right]}{1-\tfrac{1}{2}} = \dfrac{3\,069}{16}$.
		\end{itemchoice}
	 }
\end{ex}

% nguồn: Ok_dungsai_lop11_Lan2
\begin{ex}%[1D2V3-6]%[ma:eb885bdf]
	Cho  cấp số nhân($u_n$) thoả mãn $\heva{& u_4+u_6=-540 \\& u_3+u_5=180.}$
	\choiceTF 
		{\True Số  hạng $u_1=2$}
		{Gọi $q$ là công bội của  cấp số nhân, thì ba số $q$; $1$; $3$ tạo thành một cấp số cộng}
		{Số $-486$ là số hạng thứ  $5$ của cấp số nhân}
		{\True Tổng của $21$ số hạng đầu cấp số  nhân đã cho bằng $5\,230\,176\,602$}
	 \loigiai{
		\begin{itemchoice}
			\itemch Gọi $q$ là công bội và $S_{21}$ là tổng của $21$ số hạng  đầu của cấp số nhân $(u_n)$.\\
			Ta có $\heva{& u_4+u_6=-540 \\& u_3+u_5=180}\Leftrightarrow \heva{& 180q = -540 \\& u_3 + u_5 = 180}\Leftrightarrow \heva{& q  = -3 \\& u_1(9+81) = 180}\Leftrightarrow \heva{&q = -3 \\& u_1 = 2}$
			\itemch $-3$;  $1$; $3$ không thể tạo thành một cấp số cộng.
			\itemch Số $-486 = 2\cdot(-3)^5$ nên số $-486$ là  số hạng thứ $6$.
			\itemch Suy ra $S_{21} = \dfrac{u_1(1-q^{21}) }{1-q} = \dfrac{2\left[1-(-3)^{21} \right]}{1-(-3) } =  \dfrac{1+3^{21}}{2}=5\,230\,176\,602$.
		\end{itemchoice}
	 }
\end{ex}

% nguồn: Ok_dungsai_lop11_Lan2
\begin{ex}%[1D2V3-6]%[ma:64efe93a]
	Cho dãy số $\left(u_n \right)$ là một cấp số nhân có $u_1=3$, $u_3=12$  và công bội của cấp số nhân đó là số âm.
	 \choiceTF
		{\True Công bội của cấp số  cộng đó là $q=-2$}
		{Số hạng thứ $25$  của cấp số nhân đó bằng $-3\cdot2^{24}$}
		{Tổng $101$ số hạng đầu của  cấp số nhân đó bằng $1-2^{101}$}
		{\True  $u_{55}=\sqrt{u_{54} u_{56}}$}
	 \loigiai{
		\begin{itemchoice}
			\itemch Ta có $u_1 = 3$, $u_3 = 12$ và công bội $q < 0$. Suy ra $u_3 = u_1 q^2 \Leftrightarrow 3\cdot q^2  = 12 \Leftrightarrow q = -2$.
			\itemch $u_{25} = u_1\cdot q^{24} = 3\cdot(-2)^{24} =  3\cdot2^{24}$.
			\itemch Tổng số hạng đầu của cấp số nhân đó là $S_{101} = \dfrac{u_1 \left(1 - q^{101}  \right) }{1 - q} = \dfrac{3\cdot \left[1 - (-2)^{101} \right]}{1 - (-2) } = 1 + 2^{101}$.
			\itemch $\sqrt{u_{54}\cdot u_{56}} = \sqrt{\dfrac{u_{55}}{q}\cdot u_{55}\cdot  q} = \sqrt{\left(u_{55} \right)^2} = \left| u_{55} \right| = u_{55} \left(u_{55} = u_1\cdot q^{54} > 0 \right)$.
		\end{itemchoice}
	} 
\end{ex}

% nguồn: Ok_dungsai_lop11_Lan2
\begin{ex}%[1D2V3-6]%[ma:642d1285]
	Cho dãy số $(u_n)$, biết $\heva{&u_1=3 \\&u_{n+1}=4u_n-1}$ (với  $n \in \mathbb{N}^*$).
	 \choiceTF
		{Số hạng thứ  năm của dãy số là $685$}
		 {\True Đặt $v_n=u_n-\dfrac{1}{3}$ thì $(v_n)$ là cấp số nhân}
		{\True Số hạng tổng  quát $u_n=\dfrac{8}{3}\cdot4^{n-1}+\dfrac{1}{3}$}
		{\True Ta có  $S_8=58\,256$}
	\loigiai{ 
		\begin{itemchoice}
			\itemch Ta  có
			$$\begin{aligned}
			& u_2=4u_1-1=4\cdot3-1=12-1=11, \\ 
			 & u_3=4u_2-1=4\cdot11-1=44-1=43, \\
			& u_4=4u_3-1=4\cdot43-1=172-1=171,  \\
			& u_5=4u_4-1=4\cdot171-1=684-1=683. 
		\end{aligned}$$
		Vậy số hạng thứ năm của dãy số là  $u_5=683$.
		\itemch  Ta có $u_{n+1}=4u_n-1 \Leftrightarrow u_{n+1}-\dfrac{1}{3}=4\left(u_n-\dfrac{1}{3}\right)$.\\
		Đặt $v_n=u_n-\dfrac{1}{3}$, ta có  $v_{n+1}=4v_n$ suy ra $(v_n)$ là cấp số nhân.
		\itemch Ta có $v_{n+1}=4v_n$ suy ra $(v_n)$ là cấp  số nhân với $\heva{& q=4 \\& v_1=u_1-\dfrac{1}{3}=3-\dfrac{1}{3}=\dfrac{8}{3}.}$ \\
		 Suy ra $v_n=v_1\cdot q^{n-1}=\dfrac{8}{3}\cdot4^{n-1} \Rightarrow u_n=\dfrac{8}{3}\cdot4^{n-1}+\dfrac{1}{3}$.
		\itemch Ta có $S_8=u_1+u_2+\ldots+u_8=\dfrac{8}{3}\cdot\left(1+4+4^2+\ldots+4^7\right) +\dfrac{1}{3}\cdot8=58\,256$. 
	\end{itemchoice}
} 
\end{ex}

% nguồn: Ok_dungsai_lop11_Lan2
\begin{ex}%[1D2V3-6]%[ma:e25e56fa]
	Cho cấp số nhân $(u_n)$ có $u_2=6$ và  $u_3=18$.
	 \choiceTF
		{\True  Công bội của cấp số nhân là $q=3$}
		{$u_{k-1}+u_{k+1}=u_k^2$, $2 \le k \in \mathbb{N}^*$} 
		{Số hạng đầu của cấp số nhân  là $u_1=1$}
		{\True Tổng $10$ số hạng đầu tiên của cấp số nhân đã  cho $S_{10}=59048$}
	\loigiai{ 
		\begin{itemchoice}
			\itemch  Công bội của cấp số nhân là $q=\dfrac{u_3}{u_2}=3$.
			\itemch Do  $(u_n)$ là cấp số nhân $u_{k-1}.u_{k+1}=u_k^2$, $2 \le k \in \mathbb{N}^*$.
			\itemch Số hạng đầu của  cấp số nhân là $u_1=\dfrac{u_2}{q}=2$.
			\itemch Tổng $10$ số hạng đầu tiên của cấp số nhân  đã cho $S_{10}=2\cdot\dfrac{1-(3)^{10}}{1-3}=59\,048$.
		\end{itemchoice}
	} 
\end{ex}

% nguồn: Ok_dungsai_lop11_Lan1
\begin{ex}%[1D2B3-6]%[ma:5923774a]
  Hãy cho biết tính đúng sai của mỗi mệnh đề dưới đây?
 \choiceTF 
  {Tổng 10 số hạng  đầu tiên của cấp số nhân $(u_n)$, $u_1=-3$ và công bội $q=-2$ bằng $1024$}
  {\True Cho dãy số $(u_n)$  là cấp số nhân có số hạng đầu $u_1=1$, công bội $q=2$. Tổng ba số hạng đầu của cấp số nhân bằng $7$}
  {Một  cấp số nhân $(u_n)$ có $u_1=2$, $u_2=-2$. Tổng của $9$ số hạng đầu của cấp số nhân đó bằng $1$}
  {Cho  $S=11+101+1001+\cdots +\underbrace{1000\ldots 01}_{(n-1)\text{ chữ số 0}}$. Khi đó $S=10\left(\dfrac{10^n-1}{9}\right)+n$}
  \loigiai{
  \begin{itemchoice}
   \itemch  Ta có $S_{10}=\dfrac{u_1\left(1-q^{10}\right)}{1-q}=1023$.
   \itemch Tổng ba số hạng đầu của cấp số nhân có  $u_1=1$ và $q=2$ là
   $S_3=\dfrac{u_1\left(1-q^3\right)}{1-q}=\dfrac{1  \cdot\left(1-2^3\right)}{1-2}=7$.
   \itemch Công bội của cấp số  nhân đã cho là: $q=\dfrac{u_2}{u_1}=\dfrac{-2}{2}=-1$.\\
   Tổng của $9$ số hạng đầu của cấp  số nhân đó là: $S_9=2\cdot \dfrac{1-(-1)^9}{1-(-1)}=2$.
   \itemch Ta  có
   \begin{align*}
    S&=(10+1)+(10^2+1)+(10^3+1)+\cdots +(10^n+1)\\ 
    &=\left(10+10^2+10^3+\cdots  +10^n\right)+\underbrace{1+1+1+\cdots+1}_{n\text{ số } 1}\\
     &=10\left(\dfrac{10^n-1}{9}\right)+n.
   \end{align*}
  \end{itemchoice}
 } 
\end{ex}

% nguồn: Ok_dungsai_lop11_Lan1
\begin{ex}%[1D2B3-6]%[ma:9bd11965]
  Hãy cho biết tính đúng sai của mỗi mệnh đề dưới đây?
  \choiceTF
  {Cho cấp số nhân $(u_n)$ có $u_1=4, u_4=32$. Tổng $10$ số  hạng đầu tiên của cấp số nhân là $2048$}
  {\True Cho cấp số nhân $(u_n)$ có số hạng đầu $u_{1}=5$, công bội $q=-\dfrac{1}{3}$. Tổng 5 số hạng đầu của cấp số nhân đó bằng $\dfrac{305}{81}$} 
  {\True  Cho cấp số nhân $(u_n)$ với $u_1=2$, $q=4$. Tổng của $5$ số hạng đầu tiên bằng $682$}
  {Cho cấp số nhân $\left(u_n\right)$ có $u_2=-6$; $u_6=-486$. Biết $u_3>0$. Khi đó giá trị của tổng  $S=u_4+u_5+\cdots+u_{12}$ bằng $-265720$}
  \loigiai{
  \begin{itemchoice}
   \itemch Do $ u_4 = u_1 q^3 \Rightarrow q^3 = 8 \Rightarrow q = 2  $.\\
   Vậy $ S_{10} = u_1 \dfrac{q^{10} - 1}{q-1}  = 4\dfrac{2^{10}-1}{2-1} = 4092 $.
   \itemch  Ta có $S_5=u_1\cdot\dfrac{q^5-1}{q-1}=5\cdot\dfrac{\left(-\dfrac{1}{3}\right)^5-1}{\left(-\dfrac{1}{3}\right)-1}=\dfrac{305}{81}$.
   \itemch Ta có  $S_5=\dfrac{u_1(1-q^5)}{1-q}=682$.
   \itemch Ta  có $\dfrac{u_6}{u_2}=81=\dfrac{u_1\cdot q^5}{u_1\cdot q}=q^4\Rightarrow q=\pm 3$.\\
    Vì $u_2<0$, $u_3>0\Rightarrow q=-3$. Ta cũng suy ra $u_5=162$, $u_4=-54$.\\
   Vậy $S=u_4+u_5+\cdots+u_{12}=u_4\cdot\dfrac{q^9-1}{q-1}=-54\cdot  \dfrac{(-3)^9-1}{-3-1}=-265734$.
  \end{itemchoice}
 } 
\end{ex}


%%%%%%%%%%%%%%%% TRẢ LỜI NGẮN %%%%%%%%%%%%%%%%

% nguồn: Ok_TraLoiNgan_Lop11_Lan1
\begin{ex}%[1D2H3-6]%[ma:84b0c16b]
	Kết quả bốn chữa số sau cùng của  tổng sau $S_n=8+88+888+\ldots+\underbrace{88 \ldots 8}_{n \, \text {chữ số}}$ với $n=9$ bằng
	\shortans{4312} 
	 \loigiai{
		Ta  có
		$$ 
		\begin{aligned}
			S_n & =\dfrac{8}{9}(9+99+999+\ldots+\underbrace{99 \ldots 9}_{n  \, \text {chữ số}})\\
			 & =\dfrac{8}{9}\left(10-1+10^2-1+10^3-1+\ldots+10^n-1\right)\\
			& =\dfrac{8}{9}\left[\left(10+10^2+10^3+\ldots+10^n\right)-n\right]=\dfrac{8}{9}\left(10 \cdot \dfrac{1-10^n}{1-10}-n\right)=\dfrac{80\left(10^n-1\right)}{81}-\dfrac{8}{9} n. 
		\end{aligned}
		$$ 
		 Với $n=9$, ta có tổng $80\cdot\dfrac{10^9-1}{81} -\dfrac{8}{9}\cdot 9=987654312$.\\
		Vậy bốn  chữ số sau cùng của tổng trên là $4312$.
	} 
\end{ex}

% nguồn: Ok_TraLoiNgan_Lop11_Lan1
\begin{ex}%[1D2H3-6]%[ma:3ef40a20]
	Cho cấp số nhân  có các số hạng lần lượt là ${\dfrac{1}{4};\, \dfrac{1}{2};\, 1; \ldots;\, 128}$. Tính tổng ${S}$ của tất cả các số hạng của cấp số nhân đã cho (kết quả lấy phần nguyên).
	 \shortans{255}
	\loigiai{ 
		Cấp số nhân đã cho có $\heva{&u_1=\dfrac{1}{4} \\ &q=2} \Rightarrow 128=2^7=u_1 q^{n-1}=\dfrac{1}{4}  \cdot 2^{n-1}=2^{n-3} \Leftrightarrow n=10$.\\
		Khi đó  cấp số nhân đã cho có tất cả $10$ số hạng.\\
		 Suy ra $S_{10}=u_1 \cdot \dfrac{1-q^{10}}{1-q}=\dfrac{1}{4} \cdot \dfrac{1-2^{10}}{1-2}=\dfrac{1023}{4}= 255{,}75$.\\
		Vậy  $\left[S_{10} \right]=255$.
	} 
\end{ex}

% nguồn: Ok_TraLoiNgan_Lop11_Lan1
\begin{ex}%[1D2H3-6]%[ma:95159762]
	Cho cấp  số nhân $\left(u_n\right)$ có $u_1=4$ và công bội $q=-3$. Tổng của $7$ số hạng đầu tiên là
	\shortans{2,19} 
	\loigiai{ 
		Ta có  $u_1=4$, $q=-3$, và $n=7$.
		\[S_7=4\cdot\dfrac{1-(-3)^7}{1-(-3)}=2188.\] 
		Vậy, tổng của $7$ số hạng  đầu tiên là $2\,188$.
	 }
\end{ex}

% nguồn: Ok_TraLoiNgan_Lop11_Lan1
\begin{ex}%[1D2H3-6]%[ma:04858860]
	Cho cấp số nhân $(u_n)$ biết $u_1=5, u_5=405$ và  tổng $S_n=u_1+u_2+\ldots+u_n=1820$. Tìm $n$.
	\shortans{6} 
	\loigiai{ 
		Ta có: $\heva{&u_1=5 \\& u_5=405} \Leftrightarrow\heva{&u_1=5  \\& u_1 q^4=405} \Leftrightarrow\heva{&u_1=5 \\& q^4=\dfrac{405}{u_1}=81} \Leftrightarrow\heva{&u_1=5 \\& q= \pm 3.}$\\
		\textbf{Trường  hợp 1:} $u_1=5; q=3$.
		 \allowdisplaybreaks
		\begin{eqnarray*}
			&&S_n=u_1+u_2+\ldots .+u_n=1820\left(n \in  \mathbb{N}^*\right) \\
			& \Leftrightarrow&  u_1 \dfrac{1-q^n}{1-q}=1820\\
			&\Leftrightarrow&  \dfrac{1-3^n}{1-3}=\dfrac{1820}{5}\\
			&\Leftrightarrow &3^n=729\\ 
			& \Leftrightarrow&  n=6.
		\end{eqnarray*}
		\textbf{Trường hợp 2:} $u_1=5;  q=-3$.
		 \allowdisplaybreaks
		\begin{eqnarray*}
			&& S_n=u_1+u_2+\ldots  +u_n=1820\left(n\in \mathbb{N}^*\right)\\
			 &\Leftrightarrow& u_1\dfrac{1-q^n}{1-q}=1820\\
			 &\Leftrightarrow& \dfrac{1-(-3)^n}{1+3}=\dfrac{1820}{5} \\
			& \Leftrightarrow& (-3)^n=-1455\\ 
			&\Leftrightarrow& n\in \varnothing ,\forall  n\in \mathbb{N}^*.
		\end{eqnarray*}
		Vậy $n=6$. 
	 }
\end{ex}

% nguồn: Ok_TraLoiNgan_Lop11_Lan2
\begin{ex}%[1D2H3-6]%[ma:9c369d39]
	Ký hiệu $S_n$ là tổng $n$ số hạng đầu của một cấp số nhân. Biết $S_3=1$, $S_6=9$. Tính giá trị $S_{10}$ (làm tròn kết quả  đến hàng đơn vị).
	 \shortans{146}
	\loigiai{ 
		Nếu $q=1$ suy ra  $S_3=3 u_1 \Leftrightarrow u_1=\dfrac{1}{3}$. Khi đó $S_6=6 \cdot u_1=2 \neq 9$.\\
		Nếu $q \neq 1$ ta có $S_3=u \dfrac{1-q^3}{1-q}=1$.\\ 
		Lại có $S_6=u_1 \cdot \dfrac{1-q^6}{1-q}=u_1 \cdot\left(\dfrac{1-q^3}{1-q}\right) \cdot\left(1+q^3\right)=S_3 \cdot\left(1+q^3\right)=9 \Leftrightarrow q^3=8 \Leftrightarrow q=2$.\\ 
		Thay vào  ta được $u_1=\dfrac{1}{7}$.\\
		 Khi đó $S_{10}=u_1 \cdot \dfrac{1-q^{10}}{1-q}=\dfrac{1}{7}\cdot\left(\dfrac{1-2^{10}}{1-2}\right)=\dfrac{1\,023}{7} \approx 146$
	 }
\end{ex}

% nguồn: Ok_TraLoiNgan_Lop11_Lan2
\begin{ex}%[1D2H3-6]%[ma:8579e435]
	Cho cấp số  nhân $\left(u_n\right)$ thỏa mãn $\heva{&u_4-u_2=36 \\& u_3-u_3=72}$. Tính tổng $10$ số hạng đầu của cấp số nhân trên.
	 \shortans{6,32}
	 \loigiai{
		Ta  có $\heva{&36=u_4-u_2=u_1q(q^2-1) \\&
		 72=u_5-u_3=u_1q^2(q^2-1)=[u_ 1q(q^2-1)]q=36q} \Leftrightarrow \heva{&
		 q=2 \\&
		u_1=\dfrac{36}{q\left(q^2-1\right)}=6.}$\\ 
		Vậy  $ S_{10}=\dfrac{u_1\left(1-q^{10}\right)}{1-q}=\dfrac{6\left(1-2^{10}\right)}{1-2}=6\,138.$
	} 
\end{ex}

% nguồn: Ok_TraLoiNgan_Lop11_Lan1
\begin{ex}%[1D2V3-6]%[ma:5996d675]
	Cho cấp số nhân$\left(u_n\right)$ có $ \heva{&u_3=16\\&u_2+u_4=40}$. Tính  tổng $8$ số hạng đầu của cấp số nhân $\left(u_n\right)$, biết $ q > 1$.
	\shortans{1020} 
	 \loigiai{
		Ta có  $\heva{&u_3=16\\&u_2+u_4=40}\Leftrightarrow\heva{&u_1\cdot q^2=16\\&u_1\cdot q+u_1\cdot q^3=40.}$\\
		Suy ra  $\dfrac{q}{1+q^2}=\dfrac{16}{40}\Leftrightarrow 2q^2-5q+2=0\Leftrightarrow\hoac{&q=2\\&q=\dfrac{1}{2}.}$\\
		Vì $ q > 1$ nên  $ q=2$. Khi đó $u_1=4$.\\
		Tổng 8 số hạng đầu của cấp số nhân $\left(u_n\right)$ là $S_8=\dfrac{u_1\left(1-q^8\right)}{1-q}=\dfrac{4\left(1-2^8\right)}{1-2}=1020$. 
	 }
\end{ex}

% nguồn: Ok_TraLoiNgan_Lop11_Lan1
\begin{ex}%[1D2V3-6]%[ma:bb5224f4]
	Viết thêm bốn số vào giữa hai số $160$ và $5$ để được một cấp số nhân gồm sáu số hạng. Tìm tổng  tất cả các số hạng của cấp số nhân đó.
	 \shortans{315}
	\loigiai{ 
		Gọi $(u_n)$ là cấp số nhân lập được và ${q}$ là công bội của  cấp số nhân đó.\\
		Cấp số nhân cần lập  có dạng: $160; u_2; u_3; u_4; u_5; 5$.\\
		Ta có : $\heva{&u_1=160 \\  &u_6=5} \Leftrightarrow\heva{&u_1=160 \\ &u_1 q^5=5} \Leftrightarrow\heva{&u_1=160 \\ &160 q^5=5} \Leftrightarrow\heva{&u_1=160 \\& q=\dfrac{1}{2}.}$\\
		Tổng các số hạng của  cấp số nhân là: $S_6=\dfrac{u_1\left(1-q^6\right)}{1-q}=\dfrac{160\left[1-\left(\dfrac{1}{2}\right)^6\right]}{\dfrac{1}{2}}=315$.
	} 
\end{ex}

% nguồn: Ok_TraLoiNgan_Lop11_Lan1
\begin{ex}%[1D2V3-6]%[ma:61dc4f3b]
	Ông Sơn trồng cây trên một mãnh đất hình tam giác theo quy  luật: ở hàng thứ nhất có $1$ cây, ở hàng thứ hai có $2$ cây, ở hàng thứ ba có $3$ cây,..., ở hàng thứ $n$ có $n$ cây. Biết rằng ông đã trồng hết $11325$ cây. Hỏi số hàng cây được trồng theo cách trên là bao nhiêu?
	\shortans{150} 
	\loigiai{ 
		Đặt $u_1=1$, $u_2=2$, $u_3=3$,  $\ldots$, $u_n=n$.\\
		 Dãy số $\left(u_n\right)$ là một cấp số cộng với công sai $d=1$.\\
		 Ta có
		\begin{eqnarray*}
			&  & S_n=11325\\
			&\Leftrightarrow & S_n=n u_1+\dfrac{n(n-1)}{2}  \cdot d\\
			&\Leftrightarrow & 11325=n \cdot  1+\dfrac{n(n-1)}{2} \cdot 1\\
			 &\Leftrightarrow & n^2+n-22650=0\\
			 &\Leftrightarrow & \hoac{&n=150\\&n=-151.}
		\end{eqnarray*}
		Vậy  có $150$ hàng cây.
	 }
\end{ex}

% nguồn: Ok_TraLoiNgan_Lop11_Lan1
\begin{ex}%[1D2V3-6]%[ma:b0fea956]
	Tính tổng tất  cả các số hạng của một cấp số nhân, biết số hạng đầu bằng $18$, số hạng thứ hai bằng $54$ và số hạng cuối bằng $1458$.
	 \shortans{2178}
	 \loigiai{
		Gọi ${(u_n)}$ là cấp số nhân cần tìm, ${q}$ là công  bội của cấp số nhân đó.\\
		Ta có:  ${u_1=18, u_2=54 \Rightarrow q=\dfrac{u_2}{u_1}=\dfrac{54}{18}=3}$.\\
		Xét  số hạng cuối ${u_n=1458}$
		${\Leftrightarrow u_1 q^{n-1}=39366 \Leftrightarrow 18.3^{n-1}=1458 \Leftrightarrow 3^{n-1}=3^4 \Leftrightarrow n=5. 
		 }$\\
		Vậy tổng tám  số hạng của cấp số nhân là: $S_8=\dfrac{u_1\left(1-q^5\right)}{1-q}=18 \cdot \dfrac{1-3^5}{1-3}=2178$.
	 }
\end{ex}

% nguồn: Ok_TraLoiNgan_Lop11_Lan2
\begin{ex}%[1D2V3-6]%[ma:a37c687d]
	Cho năm số $a$, $b$, $c$, $d$, $e$ tạo thành một cấp số nhân theo  thứ tự đó và các số đều khác $0$, biết $\dfrac{1}{a}+\dfrac{1}{b}+\dfrac{1}{c}+\dfrac{1}{d}+\dfrac{1}{e}=10$ và tổng của chúng bằng $40$. Tính giá trị $|S|$ với $S=abcde$.
	\shortans{32} 
	\loigiai{Gọi $q (q \ne 0)$ là công bội của cấp số nhân $a$, $b$, $c$, $d$,  $e$. \\
		Khi đó $\dfrac{1}{a}$, $\dfrac{1}{b}$, $\dfrac{1}{c}$,  $\dfrac{1}{d}$, $\dfrac{1}{e}$ là cấp số nhân có công bội $\dfrac{1}{q}$.\\
		 Theo đề bài ta có
		$$\heva{&a+b+c+d+e=40 \\& \dfrac{1}{a} + \dfrac{1}{b} + \dfrac{1}{c} +  \dfrac{1}{d} + \dfrac{1}{e} = 10 }\Leftrightarrow \heva{&a\cdot\dfrac{1-q^5}{1-q} = 40 \\& \dfrac{1}{a}\cdot\dfrac{1-\left(\tfrac{1}{q}\right)^5}{1-\tfrac{1}{q}}= 10} \Leftrightarrow \heva{& a\cdot\dfrac{1-q^5}{1-q} = 40 \\& \dfrac{1}{a}\cdot\dfrac{q^5-1}{q^4(q-1) } = 10} \Rightarrow a^2q^4 = 4.$$
		Ta có $S  = abcde = a\cdot aq\cdot aq^2\cdot aq^3\cdot aq^4 = a^5q^{10}$ nên $S^2 = \left(a^5q^{10} \right)^2 = \left(a^2q^4 \right)^5$.\\
	Suy ra $\left| S \right| =  \sqrt{4^5} = 32$.}
\end{ex}

% nguồn: Ok_TraLoiNgan_Lop11_Lan2
\begin{ex}%[1D2V3-6]%[ma:86ec9391]
	Cho dãy số $\left(u_n \right)$ thỏa mãn $\heva{&5u_1+\sqrt{5u_1-u_2}=u_2+6 \\& u_{n+1}=3u_n \forall n \in \mathbb{N}^*}$. Giá trị nhỏ nhất của $n$ để $u_n \geq 2\cdot3^{2018}$ bằng bao nhiêu? 
	\shortans{2019} 
	 \loigiai{Ta có $\heva{& 5u_1 + \sqrt{5u_1 - u_2} = u_2 + 6\,(1) \\& u_{n+1} = 3u_n \forall n \in \mathbb{N}^* \,(2).}$\\
		Từ $(1)$ có $5u_1 + \sqrt{5u_1 - u_2} = u_2 + 6 \Leftrightarrow \left(\sqrt{5u_1 - u_2} \right)^2 + \sqrt{5u_1  - u_2} - 6 = 0 \Leftrightarrow 5u_1 - u_2 = 4$.\\
		Từ $(2)$ có $u_{n+1} = 3u_n \Rightarrow u_2 = 3u_1$. Giải  hệ $\heva{& 5u_1 - u_2 = 4 \\& u_2 = 3u_1}$ được $u_1 = 2$.\\
		Dãy $\left(u_n \right) $ là cấp số nhân  với $\heva{& u_1 = 2 \\& q = 3}$ có số hạng tổng quát $u_n = 2\cdot 3^{n-1}$ với $n \in \mathbb{N}^*$.\\
		Ta có $u_n \geq 2\cdot 3^{2018} \Leftrightarrow 2\cdot 3^{n-1} \geq 2\cdot 3^{2018}  \Leftrightarrow n-1 \geq 2018 \Leftrightarrow n \geq 2019$.\\
		Vậy giá  trị nhỏ nhất thỏa mãn là $2019$.
	 }
\end{ex}

% nguồn: Ok_TraLoiNgan_Lop11_Lan2
\begin{ex}%[1D2V3-6]%[ma:39a44039]
	Cho cấp số nhân $\left(u_n\right)$ thỏa $\heva{&u_1+u_2+u_3=13 \\& u_4-u_1=26}$. Tính tổng $8$ số hạng  đầu của cấp số nhân $\left(u_n\right)$.
	\shortans{3,28} 
	 \loigiai{
		Gọi $q$ là công  bội của cấp số nhân, ta có
		\[\heva{&u_1 + u_2+  u_3=13 \\&
		u_4-u_1=  26}\Leftrightarrow \heva{&
		 u_1\left(1+q+q^2\right)=13 \\&
		u_1\left(q^3-1\right)=26} \Rightarrow \dfrac{q^3-1}{1+q+q^2}=\dfrac{26}{13} \Rightarrow q-1=2  \Rightarrow q=3 \Rightarrow u_1=1.\]
		 Vậy $S_8=\dfrac{1\left(1-3^8\right)}{1-3}=3\,280$.
	} 
\end{ex}


%%%%%%%%%%%%%%%% TỰ LUẬN %%%%%%%%%%%%%%%%

% nguồn: Ok_TuanLuan_Toan11_Lan2
\begin{ex}%[1D2H3-6]%[ma:09ddf02a]
	Cho một cấp số nhân $(u_n)$ có các số  hạng đều không âm và thỏa mãn $u_2 = 6$ và $u_4 = 24$. Tính tổng của $12$ số hạng đầu tiên của cấp số nhân đó.
	\loigiai{ 
		Gọi công bội của cấp số nhân là $q$. Ta có: $u_4 = u_2  \cdot q^2$ hay $24 = 6 \cdot q^2$ suy ra $q = -2$ hoặc $q = 2$.\\
		Do cấp số nhân có các số hạng không âm  nên $q = 2$.\\
		Khi đó cấp số  nhân $(u_n)$ có $u_1 = \dfrac{u_2}{q} = \dfrac{6}{2} = 3$, suy ra
		 \[
		S_{12} = \dfrac{u_1 \left(1 - q^{12}\right)}{1 - q} =  \dfrac{3\left(1 - 2^{12}\right)}{1 - 2} = 12285.
		 \]
	} 
\end{ex}

% nguồn: Ok_TuanLuan_Toan11_Lan2
\begin{ex}%[1D2H3-6]%[ma:0e61ceef]
	 Cho cấp số nhân $\left(u_n\right)$ thỏa
	 $\heva{&u_4=\dfrac{2}{27} \\& u_3=243 u_8.}$
	Tính tổng  $10$ số hạng đầu của cấp số nhân.
	 \loigiai{
		Gọi $q$ là  công bội của cấp số. Theo giả thiết ta có
		 \[\heva{&u_1 q^3=\dfrac {2}{2 7}\\&u_1q^2=243u_1q^7} \Leftrightarrow \heva{&u_1 q^3=\dfrac {2}{2 7}\\&q ^ {5}=\dfrac {1}{2 4 3}} \Leftrightarrow \heva{& q=\dfrac {1}{3}\\&u_1=2.}
		\] 
		Tổng $10$  số hạng đầu của cấp số nhân
		\[ 
		S_{10}=u_1 \dfrac{q^{10}-1}{q-1}=2  \cdot \dfrac{\left(\dfrac{1}{3}\right)^{10}-1}{\dfrac{1}{3}-1}=3\left[1-\left(\dfrac{1}{3}\right)^{10}\right]=\dfrac{59048}{19683}.
		 \]
	} 
\end{ex}

% nguồn: Ok_TuanLuan_Toan11_Lan2
\begin{ex}%[1D2H3-6]%[ma:035d31ed]
	Tính tổng tất cả các số hạng của một cấp số nhân, biết số hạng đầu bằng $18$, số hạng thứ hai bằng $54$ và  số hạng cuối bằng $39366$.
	 \loigiai{
		Ta  có $u_1=18$, $u_2=54 \Rightarrow q=3$.\\
		$u_n=39366 \Leftrightarrow u_1 \cdot q^{n-1}=39366 \Leftrightarrow  18 \cdot 3^{n-1}=39366 \Leftrightarrow 3^{n-1}=3^7 \Leftrightarrow n=8.$\\
		Vậy $S_8=18 \cdot \dfrac{1-3^8}{1-3}=59040$. 
	} 
\end{ex}

% nguồn: Ok_TuanLuan_Toan11_Lan2
\begin{ex}%[1D2V3-6]%[ma:9dd4f2a7]
	Tính các tổng sau: 
	\begin{enumerate}
		\item $S_n = 2 + 2^2 + 2^3 + \cdots + 2^n$. 
		 \item $S_n = \dfrac{1}{2} + \dfrac{1}{2^2} + \dfrac{1}{2^3} + \cdots + \dfrac{1}{2^n}$.
		\item $S_n = \left(3 +  \dfrac{1}{3}\right)^2 + \left(9 + \dfrac{1}{9}\right)^2 + \cdots + \left(3^n + \dfrac{1}{3^n}\right)^2$.
		\item  $S_n = 6 + 66 + 666 + \cdots + \underbrace{666\ldots6}_{n\text{ số }6}$.
	\end{enumerate}
	\loigiai{ 
		\begin{enumerate}
			\item Dãy số $2$, $2^2$, $2^3$, $\cdots$, $2^n$ là một cấp số nhân với $n$ số hạng, có số hạng đầu $u_1 = 2$ và công bội $q = \dfrac{2^2}{2} = 2$.  Do đó
			\[ 
			S_n  = u_1 \cdot \dfrac{1 - q^n}{1 - q} = 2 \cdot \dfrac{1 - 2^n}{1 - 2} = 2(2^n - 1).
			\] 
			\item Dãy số $\dfrac{1}{2}$, $\dfrac{1}{2^2}$, $\dfrac{1}{2^3}$, $\cdots$, $\dfrac{1}{2^n}$ là một cấp số nhân với $n$ số hạng, có số  hạng đầu $u_1 = \dfrac{1}{2}$ và công bội $q = \dfrac{\tfrac{1}{2^2}}{\tfrac{1}{2}} = \dfrac{1}{2}$. Do đó
			\[ 
			S_n = u_1 \cdot \dfrac{1 - q^n}{1 - q}  = \dfrac{1}{2} \cdot \dfrac{1 - \left(\dfrac{1}{2}\right)^n}{1 - \dfrac{1}{2}} = 1 - \dfrac{1}{2^n}.
			\] 
			\item $S_n = \left(3 + \dfrac{1}{3}\right)^2 + \left(9 + \dfrac{1}{9}\right)^2 + \cdots +  \left(3^n + \dfrac{1}{3^n}\right)^2$\\
			$= (3^2 + 3^4 + \cdots + 3^{2n}) +  \left(\dfrac{1}{3^2} + \dfrac{1}{3^4} + \cdots + \dfrac{1}{3^{2n}}\right) + \underbrace{2 + 2 + \cdots + 2}_{n\text{ số }2}$.\\
			 Ta có dãy số $3^2$, $3^4$, $\cdots$, $3^{2n}$ là một cấp số nhân với $n$ số hạng, $u_1 = 3^2$, $q = \dfrac{3^4}{3^2} = 9$.\\
			 Do đó:
			\[ 
			S_1 = u_1 \cdot \dfrac{1 - q^n}{1 -  q} = 9 \cdot \dfrac{1 - 9^n}{1 - 9} = \dfrac{9}{8}(9^n - 1).
			 \]
			Tương tự, dãy $\dfrac{1}{3^2}, \dfrac{1}{3^4}, \cdots, \dfrac{1}{3^{2n}}$ là cấp số nhân với $u_1 =  \dfrac{1}{3^2}$, $q = \dfrac{1}{9}$. Do đó:
			\[ 
			 S_2 = u_1 \cdot \dfrac{1 - q^n}{1 - q} = \dfrac{1}{9} \cdot \dfrac{1 - \left(\dfrac{1}{9}\right)^n}{1 - \dfrac{1}{9}} = \dfrac{1}{8} \left(1 - \dfrac{1}{9^n} \right) = \dfrac{9^n - 1}{8 \cdot 9^n}.
			\] 
			Vậy: 
			 \[
			S_n = \dfrac{9}{8}(9^n - 1) + \dfrac{9^n - 1}{8 \cdot 9^n} + 2n = \dfrac{(9^n  - 1)(9^{n+1} + 1)}{8 \cdot 9^n} + 2n.
			\] 
			 \item
			\begin{align*}
				 \allowdisplaybreaks
				S_n &= 6 +  66 + 666 + \cdots + \underbrace{666\ldots6}_{n\text{ số }6} \\
				&= \dfrac{6}{9} \left(9 + 99 + 999  + \cdots + \underbrace{999\ldots9}_{n\text{ số }9}\right) \\
				&= \dfrac{2}{3}\big[(10 - 1) + (100  - 1) + (1000 - 1) + \cdots + (10^n - 1)\big]\\
				&= \dfrac{2}{3} \left[10 +  10^2 + 10^3 + \cdots + 10^n - n\right]\\
				&= \dfrac{2}{3} \left[\dfrac{10^{n+1} -  10}{9} - n\right]\\
				&= \dfrac{20}{27}(10^n - 1) - \dfrac{2n}{3}. 
			\end{align*}
		\end{enumerate}
	 }
\end{ex}

% nguồn: Ok_TuanLuan_Toan11_Lan2
\begin{ex}%[1D2V3-6]%[ma:21937cb8]
	Tính tổng sau  $S_n=\left(2+\dfrac{1}{2}\right)^2+\left(4+\dfrac{1}{4}\right)^2+\cdots+\left(2^n+\dfrac{1}{2^n}\right)^2$.
	 \loigiai{
		 \[
		\begin{aligned}
			S_n&= 2^2+\dfrac{1}{2^2}+2+2^4+\dfrac{1}{2^4}+2+\cdots+2^{2 n}+\dfrac{1}{2^{2 n}}+2  \\
			&=\left(2^2+2^4+\cdots+2^{2 n}\right)+\left(\dfrac{1}{2^2}+\dfrac{1}{2^4}+\cdots+\dfrac{1}{2^{2  n}}\right)+2 n \\
			 &=4 \cdot \dfrac{1-4^n}{1-4}+\dfrac{1}{4}\cdot \dfrac{1-\dfrac{1}{4^n}}{1-\dfrac{1}{4}}+2 n=\dfrac{4^n-1}{3}\left(4-\dfrac{1}{4^n}\right)+2n.
		\end{aligned}
		 \]
	 }
\end{ex}

% nguồn: Ok_TuanLuan_Toan11_Lan2
\begin{ex}%[1D2V3-6]%[ma:a25ff6c8]
	 $S_n=8+88+888+\cdots+\underbrace{88 \ldots 8}_{n\, \text{số}\, 8}$.
	\loigiai{ 
		 \[
		\begin{aligned}
			S_n&= \dfrac{8}{9}(9+99+999+\underbrace{99 \ldots 9}_{n\, \text{số}\,  9}) \\
			&=\dfrac{8}{9}\left(10-1+10^2-1+10^3-1+\cdots+10^n-1\right)  \\&=\dfrac{8}{9}{\left[\left(10+10^2+10^3+\cdots+10^n\right)-n\right]} \\
			&=\dfrac{8}{9}\left(10 \cdot \dfrac{1-10^n}{1-10}-n\right)=\dfrac{80\left(10^n-1\right)}{81}-\dfrac{8}{9n}. 
		\end{aligned}
		\] 
	} 
\end{ex}
